\section{JAM 2005}
\begin{enumerate}
\item 
 Let $\left\{a_{n}\right\},\left\{b_{n}\right\}$ and $\left\{c_{n}\right\}$ be sequences of real numbers such that $b_{n}=a_{2 n}$ and $c_{n}=a_{2 n+1}$. Then $\left\{a_{n}\right\}$ is convergent
 \begin{enumerate}
     \item implies $\left\{b_{n}\right\}$ is convergent but $\left\{c_{n}\right\}$ need not be convergent

     \item implies $\left\{c_{n}\right\}$ is convergent but $\left\{b_{n}\right\}$ neet not be convergent

     \item implies both $\left\{b_{n}\right\}$ and $\left\{c_{n}\right\}$ are convergent

     \item if both $\left\{b_{n}\right\}$ and $\left\{c_{n}\right\}$ are convergent
 \end{enumerate}
 \hfill [JAM 2005]


\item  An integrating factor of $x \frac{d y}{d x}+(3 x+1) y=x e^{-2 x}$ is
 \begin{enumerate}
     \item  $x e^{3 x}$
 \item  $3 x e^{x}$
 \item  $x e^{x}$
 \item  $x^{3} e^{x}$ 
 \end{enumerate}
 \hfill [JAM 2005]
 
\item The general solution of $x^{2} \frac{d^{2} y}{d x^{2}}-5 x \frac{d y}{d x}+9 y=0$ is
\begin{enumerate}
     \item  $\left(c_{1}+c_{2} x\right) e^{3 x}$
     \item  $\left(c_{1}+c_{2} \ln x\right) x^{3}$
     \item  $\left(c_{1}+c_{2} x\right) x^{3}$
     \item $\left(c_{1}+c_{2} \ln x\right) e^{x^{3}}$
 \end{enumerate}
(Here $c_{1}$ and $c_{2}$ are arbitrary constants.)
 \hfill [JAM 2005]




\item  Let $\vec{r}=x \hat{i}+y \hat{j}+z \hat{k}$. If $\phi(x, y, z)$ is a solution of the Laplace equation then the vector field $(\vec{\nabla} \phi+\vec{r})$ is
\begin{enumerate}
    \item  neither solenoidal nor irrotational
    \item solenoidal but not irrotational
    \item both solenoidal and irrotational
    \item irrotational but not solenoidal
 \end{enumerate}
 \hfill [JAM 2005]



 
\item  Let $\vec{F}=x \hat{i}+2 y \hat{j}+3 z \hat{\mathrm{k}}, S$ be the surface of the sphere $x^{2}+y^{2}+z^{2}=1$ and $\hat{\mathrm{n}}$ be the inward unit normal vector to $S$. Then $\iint_{S} \vec{F} \cdot \hat{n} d S$ is equal to
\begin{enumerate} 
    \item $4 \pi$
    \item $-4 \pi$
    \item $8 \pi$
    \item $-8 \pi$ 
\end{enumerate}
 \hfill [JAM 2005]



 
\item Let $A$ be a $3 \times 3$ matrix with eigenvalues $1,-1$ and 3 . Then
\begin{enumerate}
    \item $A^{2}+A$ is non-singular
    \item $A^{2}-A$ is non-singular
    \item $A^{2}+3 A$ is non-singular
    \item $A^{2}-3 A$ is non-singular
 \end{enumerate}
 \hfill [JAM 2005]

 
\item Let $T: \boldsymbol{R}^{3} \longrightarrow \boldsymbol{R}^{3}$ be a linear transformation and $I$ be the identity transformation of $\boldsymbol{R}^{3}$. If there is a scalar $c$ and a non-zero vector $x \in \boldsymbol{R}^{3}$ such that $T(x)=c x$, then $\operatorname{rank}(T-c I)$
\begin{enumerate} 
    \item cannot be 0
    \item cannot be 1
    \item cannot be 2
    \item cannot be 3
\end{enumerate}
 \hfill [JAM 2005]



 
\item  In the group $\{1,2, \ldots, 16\}$ under the operation of multiplication modulo 17, the order of the element 3 is
\begin{enumerate} 
    \item 4
    \item 8
    \item 12
    \item 16
 \end{enumerate}
 \hfill [JAM 2005]


 
\item A ring $R$ has maximal ideals
\begin{enumerate} 
    \item if $R$ is infinite
    \item if $R$ is finite
    \item if $R$ is finite with at least 2 elements
    \item only if $R$ is finite 
\end{enumerate}
 \hfill [JAM 2005]



\item  The integral 
$\int_{0}^{1}\left[\int_{0}^{1-z}\left(\int_{0}^{2} d x\right) d y\right] d z$ is equal to
\begin{enumerate}
    \item $\quad \int_{0}^{1}\left[\int_{0}^{1-y}\left(\int_{0}^{2} d x\right) d z\right] d y$
    \item $\quad \int_{0}^{1}\left[\int_{0}^{1-y}\left(\int_{0}^{1} d x\right) d z\right] d y$
    \item $\int_{0}^{2}\left[\int_{0}^{2}\left(\int_{0}^{1-z} d x\right) d z\right] d y$
    \item $\int_{0}^{2}\left[\int_{0}^{2}\left(\int_{0}^{1-y} d x\right) d z\right] d y$
 \end{enumerate}
 \hfill [JAM 2005]


 
\item Let $f: \boldsymbol{R} \longrightarrow \boldsymbol{R}$ be continuous and $g, h: \boldsymbol{R}^{2} \longrightarrow \boldsymbol{R}$ be differentiable. Let $F(u, v)=\int_{v}^{u} f(t) d t$, where $u=g(x, y)$ and $v=h(x, y)$. Then $\frac{\partial F}{\partial x}+\frac{\partial F}{\partial y}=$
\begin{enumerate} 
    \item $f(g(x, y))\left[\frac{\partial g}{\partial x}+\frac{\partial g}{\partial y}\right]-f(h(x, y))\left[\frac{\partial h}{\partial x}+\frac{\partial h}{\partial y}\right]$
    \item $f(h(x, y))\left[\frac{\partial g}{\partial x}+\frac{\partial g}{\partial y}\right]-f(g(x, y))\left[\frac{\partial h}{\partial x}-\frac{\partial h}{\partial y}\right]$
    \item $f(h(x, y))\left[\frac{\partial g}{\partial x}+\frac{\partial g}{\partial y}\right]+f(g(x, y))\left[\frac{\partial h}{\partial x}-\frac{\partial h}{\partial y}\right]$
    \item $f(g(x, y))\left[\frac{\partial g}{\partial x}-\frac{\partial g}{\partial y}\right]+f(g(x, y))\left[\frac{\partial h}{\partial x}-\frac{\partial h}{\partial y}\right]$ 
 \end{enumerate}
 \hfill [JAM 2005]



\item Let $y=f(x)$ be a smooth curve such that $0<f(x)<K$ for all $x \in[a, b]$. Let

$L=$ length of the curve between $x=a$ and $x=b$

$A=$ area bounded by the curve, $x$-axis, and the lines $x=a$ and $x=b$

$S=$ area of the surface generated by revolving the curve about $x$-axis between $x=a$ and $x=b$

Then
\begin{enumerate} 
    \item $2 \pi K L<S<2 \pi A$
    \item $S \leq 2 \pi A<2 \pi K L$
    \item $2 \pi A \leq S<2 \pi K L$
    \item $2 \pi A<2 \pi K L<S$
 \end{enumerate}
 \hfill [JAM 2005]


 
\item  Let $f: \boldsymbol{R} \longrightarrow \boldsymbol{R}$ be defined by $f(t)=t^{2}$ and let $U$ be any non-empty open subset of $\boldsymbol{R}$. Then
\begin{enumerate}
    \item $f(U)$ is open
    \item $f^{-1}(U)$ is open
    \item $f(U)$ is closed
    \item $f^{-1}(U)$ is closed 
 \end{enumerate}
 \hfill [JAM 2005]
 
 
 
 
 \item Let $f:(-1,1) \longrightarrow \boldsymbol{R}$ be such that $f^{(n)}(x)$ exists and $\left|f^{(n)}(x)\right| \leq 1$ for every $n \geq 1$ and for every $x \in(-1,1)$. Then $f$ has a convergent power series expansion in a neighbourhood of
\begin{enumerate} 
    \item every $x \in(-1,1)$
    \item every $x \in\left(\frac{-1}{2}, 0\right)$ only
    \item no $x \in(-1,1)$
    \item every $x \in\left(0, \frac{1}{2}\right)$ only
\end{enumerate}
 \hfill [JAM 2005]
 
 
 
 
 \item  Let $a>1$ and $f, g, h:[-a, a] \longrightarrow \boldsymbol{R}$ be twice differentiable functions such that for some $c$ with $0<c<1<a$,

$$
\begin{aligned}
&f(x)=0 \text { only for } x=-a, 0, a \\
&f^{\prime}(x)=0=g(x) \text { only } x=-1,0,1 \\
&g^{\prime}(x)=0=h(x) \text { only for } x=-c, c .
\end{aligned}
$$

The possible relations between $f, g, h$ are
\begin{enumerate} \item $f=g^{\prime}$ and $h=f^{\prime}$
    \item $f^{\prime}=g$ and $g^{\prime}=h$
    \item $f=-g^{\prime}$ and $h^{\prime}=g$
    \item $f=-g^{\prime}$ and $h^{\prime}=f$ 
\end{enumerate}
 \hfill [JAM 2005]



 
\begin{center}
\textbf{CORE SECTION}
\end{center}


 \item
 \begin{enumerate}
     \item Solve the initial value problem
    $$
     \frac{d^{2} y}{d x^{2}}-y=x\left(\sin x+e^{x}\right), \quad y(0)=y^{\prime}(0)=1
    $$
    
    \item Solve the differential equation

    $$
        \left(2 y \sin x+3 y^{4} \sin x \cos x\right) d x-\left(4 y^{3} \cos ^{2} x+\cos x\right) d y=0
    $$

 \end{enumerate}
  \hfill [JAM 2005]



\item  Let $G$ be a finite abelian group of order $n$ with identity $e$. If for all $a \in G, a^{3}=e$, then, by induction on $n$, show that $n=3^{k}$ for some non-negative integer $k$. 
 \hfill [JAM 2005]

\item 
\begin{enumerate}
   \item Let $f:[a, b] \longrightarrow \boldsymbol{R}$ be a differentiable function. Show that there exist points $c_{1}, c_{2} \in(a, b)$ such that

    $$
        2 f\left(c_{1}\right) f^{\prime}\left(c_{1}\right)=f^{\prime}\left(c_{2}\right)[]
    $$

    \item Let
    $$
        f(x, y)= \begin{cases}\left(x^{2}+y^{2}\right)\left[\ln \left(x^{2}+y^{2}\right)+1\right] & \text { for }(x, y) \neq(0,0) \\ \alpha & \text { for }(x, y)=(0,0)\end{cases}
    $$

    Find a suitable value for $\alpha$ such that $f$ is continuous. For this value of $\alpha$, is $f$ differentiable at $(0,0)$ ? Justify your claim. 

\end{enumerate}
 \hfill [JAM 2005]


 
\item 
\begin{enumerate}
    \item Let $S$ be the surface $x^{2}+y^{2}+z^{2}=1, z \geq 0$. Use Stokes' theorem to evaluate

        $$
            \int_{C}[(2 x-y) d x-y d y-z d z]
        $$

    where $C$ is the circle $x^{2}+y^{2}=1, z=0$, oriented anticlockwise.

    \item  Show that the vector field $\vec{F}=\left(2 x y-y^{4}+3\right) \hat{\mathrm{i}}+\left(x^{2}-4 x y^{3}\right) \hat{\mathrm{j}}$ is conservative. Find its potential and also the work done in moving a particle from $(1,0)$ to (2,1) along some curve. 

\end{enumerate}
 \hfill [JAM 2005]

 

\item  Let $T: \boldsymbol{R}^{3} \longrightarrow \boldsymbol{R}^{3}$ be defined by $T(x, y, z)=(y+z, z, 0)$. Show that $T$ is a linear transformation. If $v \in \boldsymbol{R}^{3}$ is such that $T^{2}(v) \neq 0$, then show that $B=\left\{v, T(v), T^{2}(v)\right\}$ forms a basis of $\boldsymbol{R}^{3}$. Compute the matrix of $T$ with respect to $B$. Also find a $v \in \boldsymbol{R}^{3}$ such that $T^{2}(v) \neq 0$ 
 \hfill [JAM 2005]

\item 
\begin{enumerate}
    \item For each $n \in N$, define $f_{n}:[-1,1] \longrightarrow \boldsymbol{R}$ by

    $$
        f_{n}(x)= \begin{cases}4 n^{2} x & \text { for } x \in\left[0, \frac{1}{2 n}\right) \\ -4 n^{2}\left(x-\frac{1}{n}\right) & \text { for } x \in\left[\frac{1}{2 n}, \frac{1}{n}\right) \\ 0 & \text { for } x \in\left[\frac{1}{n}, 1\right]\end{cases}
    $$

    Compute $\int_{0}^{1} f_{n}(x) d x$ for each $n$. Analyse pointwise and uniform convergence of the sequence of functions $\left\{f_{n}\right\}$.
    

    \item Let $f: \boldsymbol{R} \longrightarrow \boldsymbol{R}$ be a continuous function with $|f(x)-f(y)| \geq|x-y|$ for every $x, y \in \boldsymbol{R}$. Is $f$ one-one? Show that there cannot exist three points $a, b, c \in \boldsymbol{R}$ with $a<b<c$ such that $f(a)<f(c)<f(b)$ 
\end{enumerate}
 \hfill [JAM 2005]

 
\item Find the volume of the cylinder with base as the disk of unit radius in the $x y$-plane centred at $(1,1,0)$ and the top being the surface $z=\left[(x-1)^{2}+(y-1)^{2}\right]^{3 / 2}$. 
 \hfill [JAM 2005]


 
\begin{center}
\textbf{OPTIONAL SECTION : A}
\end{center}

\item 
\begin{enumerate}
    \item Bag A contains 3 white and 4 red balls, and bag B contains 6 white and 3 red balls. A biased coin, twice as likely to come up heads as tails, is tossed once. If it shows head, a ball is drawn from bag A, otherwise, from bag B. Given that a white ball was drawn, what is the probability that the coin came up tail?

    \item Let the random variables $X$ and $Y$ have the joint probability density function $f(x, y)$ given by

        $$
            f(x, y)= \begin{cases}y^{2} e^{-y(x+1)} & x \geq 0, y \geq 0 \\ 0 & \text { otherwise }\end{cases}
        $$

    Are the random variables $X$ and $Y$ independent? Justify your answer. 
\end{enumerate}
 \hfill [JAM 2005]


\item 
\begin{enumerate}
    \item Let $X_{1}, X_{2}, \ldots, X_{n}$ be independently identically distributed random variables (rv's) with common probability density function (pdf) $f_{X}(x, \theta)=\frac{1}{\theta} e^{-x / \theta} ; x>0, \theta>0$. Obtain the moment generating function (mgf) of $\bar{X}=\frac{1}{n} \sum_{i=1}^{n} X_{i}$. Also find the mgf of the rv $Y=2 n \bar{X} / \theta$

    \item Let $X_{1}, X_{2}, \ldots, X_{9}$ be an independent random sample from $N(2,4)$ and $Y_{1}, Y_{2}, Y_{3}, Y_{4}$ be an independent random sample from $N(1,1)$. Find $P(\bar{X}>\bar{Y})$, where $\bar{X}$ and $\bar{Y}$ are sample means.

[   Given $P(Z>1.2)=0.1151$, where $Z \sim N(0,1)]$ 
\end{enumerate}
 \hfill [JAM 2005]

\begin{center}
\textbf{OPTIONAL SECTION : B}
\end{center}

\item 
\begin{enumerate}
    \item Let $X_{1}, X_{2}, \ldots, X_{n}$ be a random sample from a distribution having pdf

    $$
        f\left(x ; x_{0}, \alpha\right)= \begin{cases}\frac{\alpha x_{0}^{\alpha}}{x^{\alpha+1}} & \text { for } x>x_{0} \\ 0 & \text { otherwise }\end{cases}
    $$

    where $x_{0}>0, \alpha>0$. Find the maximum likelihood estimator of $\alpha$ if $x_{0}$ is known.

    \item Let $X_{1}, X_{2}, \ldots, X_{5}$ be a random sample from the standard normal population. Determine the constant $c$ such that the random variable

    $$
        Y=\frac{c\left(X_{1}+X_{2}\right)}{\sqrt{X_{3}^{2}+X_{4}^{2}+X_{5}^{2}}}
    $$

    will have a $t$-distribution. 
\end{enumerate}
 \hfill [JAM 2005]


 

\item 
\begin{enumerate}
    \item A random sample of size $n=1$ is drawn from pdf $f_{X}(x, \theta)=\frac{1}{\theta} e^{-x / \theta} ; x>0, \theta>0$. It is decided to test $H_{0}: \theta=5$ against $H_{1}: \theta=7$ based on the criterion: reject $H_{0}$ if the observed value is greater than $10$. Obtain the probabilities of type I and type II errors.

    \item Let $X_{1}, X_{2}, \ldots, X_{n}$ be a random sample from a normal population $N\left(\mu, \sigma^{2}\right)$. Find the best test for testing $H_{0}: \mu=0, \sigma^{2}=1$ against $H_{1}: \mu=1, \sigma^{2}=4$. 
\end{enumerate}
 \hfill [JAM 2005]


 
\begin{center}
\textbf{OPTIONAL SECTION : C}
\end{center}


\item 
\begin{enumerate}
    \item Let $f, g: \boldsymbol{R} \longrightarrow \boldsymbol{R}$ be such that for $x, y \in \boldsymbol{R}$,

    $$
        \phi(x+i y)=e^{x}[f(y)+i g(y)]
    $$

    is an analytic function. Find a differential equation of order 2 satisfied by $f$.

    \item Compute $\int_{|z+1|=2}(2 z+1) e^{(\sqrt{2}+1 / z)} d z$.  
\end{enumerate}
 \hfill [JAM 2005]


 
\item 
\begin{enumerate}
    \item  Let $f(z)$ be analytic in the whole complex plane such that for all $r>0$,

    $$
        \int_{0}^{2 \pi}\left|f\left(r e^{i \theta}\right)\right| d \theta \leq \sqrt{r}
    $$

    Find $\frac{f^{(n)}(0)}{n !}$ for all $n \geq 0$

    \item Find all values of $\alpha \in C$ such that $f(z)=(z+\bar{z})^{2}+2 \alpha|z|^{2}+\alpha(\bar{z})^{2}$ is analytic at some point $z$ having non-zero real part. 
\end{enumerate}
 \hfill [JAM 2005]

 
\begin{center}
\textbf{OPTIONAL SECTION: D}
\end{center}


\item  A hemispherical bowl of radius $12 \mathrm{~cm}$ is fixed such that its rim is horizontal. A light rod of length $20 \mathrm{~cm}$ with weights $w$ and $W$ attached to its two ends is placed inside the bowl. In equilibrium, the weight $w$ is just touching the rim of the bowl. Find the ratio $w: W$.
 \hfill [JAM 2005]




\item  A uniform ladder of length $2 a$ and mass $m$ lies in a vertical plane with one end against a smooth wall, the other end being supported on a horizontal floor. The ladder is released from rest when inclined at an angle $\alpha$ to the horizontal. Find the inclination of the ladder to the horizontal when it ceases to touch the wall.
 \hfill [JAM 2005]




\item 
\begin{enumerate}
    \item Estimate the error in evaluating the integral $\int_{0}^{8}\left(1+x^{2}\right) e^{-x} d x$ by Simpson's $\frac{1}{3}$ rd rule with spacing $h=0.25$.

    \item Using Newton-Raphson method, compute the point of intersection of the curves $y=x^{3}$ and $y=8 x+4$ near the point $x=3$, correct up to 2 decimal places. [Round-off the first iteration up to 2 decimal places for further computation]
\end{enumerate}
 \hfill [JAM 2005]



 
\item  The polynomial $p_{3}(x)=x^{3}+x^{2}-2$ interpolates the function $f(x)$ at the points $x=-1,0,1$ and 2 . If the data $f(3)=-14$ is added, find the new interpolating polynomial by using Newton's forward difference formula. Also find $f(2.5)$ by using Newton's backward difference formula with pivot value 3. Justify whether the value obtained will be the same if pivot value 2 is taken.
 \hfill [JAM 2005]


\end{enumerate}

\newpage

\section{JAM 2006}

\begin{enumerate}
  \item $\lim _{n \rightarrow \infty} \frac{2^{n+1}+3^{n+1}}{2^{n}+3^{n}}$ equals
\begin{enumerate}
    \item 3
    \item  2
    \item  1
    \item 0
\end{enumerate}
 \hfill [JAM 2006]


  \item Let $f(x)=(x-2)^{17}(x+5)^{24}$. Then
  \begin{enumerate}
    \item $f$ does not have a critical point at 2
    \item $f$ has a minimum at 2
    \item $f$ has a maximum at 2
    \item $f$ has neither a minimum nor a maximum at 2
  \end{enumerate}
 \hfill [JAM 2006]


  \item Let $f(x, y)=x^{5} y^{2} \tan ^{-1}\left(\frac{y}{x}\right)$. Then $x \frac{\partial f}{\partial x}+y \frac{\partial f}{\partial y}$ equals
    \begin{enumerate}
\item $2 f$
\item $3 f$
\item $5 f$
\item $7 f$
\end{enumerate}
 \hfill [JAM 2006]
 
  \item Let $G$ be the set of all irrational numbers. The interior and the closure of $G$ are denoted by $G^{0}$ and $\bar{G}$, respectively. Then
  \begin{enumerate}
\item $G^{0}=\phi, \quad \bar{G}=G$
\item $G^{0}=\pm, \quad \bar{G}=\pm$
\item $G^{0}=\phi, \quad \bar{G}=\bar{\Xi}$
\item $G^{0}=G, \quad \bar{G}=\ddagger$ 
\end{enumerate}
 \hfill [JAM 2006]


\item  Let $f(x)=\int_{\sin x}^{\cos x} e^{-t^{2}} d t$. Then $f^{\prime}(\pi / 4)$ equals
\begin{enumerate} 
    \item $\sqrt{1 / e}$
    \item $-\sqrt{2 / e}$
    \item $\sqrt{2 / e}$
    \item $-\sqrt{1 / e}$
\end{enumerate}
 \hfill [JAM 2006]

 

\item Let $C$ be the circle $x^{2}+y^{2}=1$ taken in the anti-clockwise sense. Then the value of the integral
$$
\int_{C}\left[\left(2 x y^{3}+y\right) d x+\left(3 x^{2} y^{2}+2 x\right) d y\right]
$$
equals
\begin{enumerate}
\item 1
\item $\pi / 2$
\item  $\pi$
\item  0
\end{enumerate}
 \hfill [JAM 2006]

  \item Let $r$ be the distance of a point $P(x, y, z)$ from the origin $O$. Then $\nabla r$ is a vector
  \begin{enumerate}
    \item orthogonal to $\overrightarrow{\mathrm{OP}}$
    \item normal to the level surface of $r$ at $P$
    \item normal to the surface of revolution generated by OP about $x$-axis
    \item normal to the surface of revolution generated by OP about y-axis
\end{enumerate}
 \hfill [JAM 2006]


\item Let $T: \dot{D}^{3} \rightarrow \dot{L}^{3}$ be defined by
$$
T\left(x_{1}, x_{2}, x_{3}\right)=\left(x_{1}-x_{2}, x_{1}-x_{2}, 0\right) .
$$
If $N(T)$ and $R(T)$ denote the null space and the range space of $T$ respectively, then
\begin{enumerate}
    \item $\operatorname{dim} N(T)=2$
    \item $\quad \operatorname{dim} R(T)=2$
    \item $R(T)=N(T)$
    \item $\quad N(T) \subset R(T)$
\end{enumerate}
 \hfill [JAM 2006]




  \item Let $\mathrm{S}$ be a closed surface for which $\iint_{S} \vec{r} . \hat{n} d \sigma=1$. Then the volume enclosed by the surface is
\begin{enumerate}
\item 1
\item $1 / 3$
\item $2 / 3$
\item 3
\end{enumerate}
 \hfill [JAM 2006]



\item  If $\left(c_{1}+c_{2} \ln x\right) / x$ is the general solution of the differential equation
$$
x^{2} \frac{d^{2} y}{d x^{2}}+k x \frac{d y}{d x}+y=0, \quad x>0,
$$

then $k$ equals
\begin{enumerate}
\item 3
\item $-3$
\item 2
\item $-1$
\end{enumerate}
 \hfill [JAM 2006]



 
  \item If $A$ and $B$ are $3 \times 3$ real matrices such that $\operatorname{rank}(A B)=1$, then $\operatorname{rank}(B A)$ cannot be
  \begin{enumerate}
\item 0
\item 1
\item 2
\item 3
\end{enumerate}
 \hfill [JAM 2006]



  \item The differential equation representing the family of circles touching y-axis at the origin is
  \begin{enumerate}
\item linear and of first order
\item linear and of second order
\item nonlinear and of first order
\item nonlinear and of second order
\end{enumerate}
 \hfill [JAM 2006]




  \item Let $\mathrm{G}$ be a group of order 7 and $\phi(x)=x^{4}, x \in G$. Then $\phi$ is

  \begin{enumerate}
\item not one - one
\item not onto
\item not a homomorphism
\item one-one, onto and a homomorphism
\end{enumerate}
 \hfill [JAM 2006]




  \item Let $\mathrm{R}$ be the ring of all $2 \times 2$ matrices with integer entries. Which of the following subsets of $R$ is an integral domain?

  \begin{enumerate}
\item $\left\{\left(\begin{array}{ll}0 & x \\ y & 0\end{array}\right): x, y \in \mathbf{Z}\right\}$

\item $\left\{\left(\begin{array}{ll}x & 0 \\ 0 & y\end{array}\right): x, y \in \mathbf{Z}\right\}$

\item $\left\{\left(\begin{array}{ll}x & 0 \\ 0 & x\end{array}\right): x \in \mathbf{Z}\right\}$

\item $\left\{\left(\begin{array}{ll}x & y \\ y & z\end{array}\right): x, y, z \in \mathbf{Z}\right\}$
\end{enumerate}
 \hfill [JAM 2006]



\item Let $f_{n}(x)=n \sin ^{2 n+1} x \cos x$. Then the value of is
$$
\lim _{n \rightarrow \infty} \int_{0}^{\pi / 2} f_{n}(x) d x-\int_{0}^{\pi / 2}\left(\lim _{n \rightarrow \infty} f_{n}(x)\right) d x
$$
\begin{enumerate}
\item $1 / 2$
\item 0
\item $-1 / 2$
\item $-\infty$
\end{enumerate}
 \hfill [JAM 2006]



\item 
\begin{enumerate}
  \item   Test the convergence of the series
$$
\sum_{n=1}^{\infty} \frac{n^{n}}{n ! 3^{n}}
$$
\item  Show that for $0 \leq x \leq \pi / 2$.
$$
\ln (1+\cos x) \leq \ln 2-\frac{x^{2}}{4}
$$
\end{enumerate}
 \hfill [JAM 2006]

 
\item Find the critical points of the function

$$
f(x, y)=x^{3}+y^{2}-12 x-6 y+40 .
$$
Test each of these for maximum and minimum.
 \hfill [JAM 2006]

 


\item 
\begin{enumerate}
  \item  Evaluate $\iint_{R} x e^{y^{2}} d x d y$, where $R$ is the region bounded by the lines $x=0, y=1$ and the parabola $y=x^{2}$.

\item Find the volume of the solid bounded above by the surface $z=1-x^{2}-y^{2}$ and below by the plane $z=0$.
\end{enumerate}
 \hfill [JAM 2006]



\item
Evaluate the surface integral
$$
\iint_{S} x\left(12 y-y^{4}+z^{2}\right) d \sigma,
$$
where the surface $S$ is represented in the form $z=y^{2}, 0 \leq x \leq 1,0 \leq y \leq 1$.
 \hfill [JAM 2006]

\item Using the change of variables, evaluate $\iint_{R} x y d x d y$, where the region $R$ is bounded by the curves $x y=1, x y=3, y=3 x$ and $y=5 x$ in the first quadrant.
 \hfill [JAM 2006]




\item 
\begin{enumerate}
\item  Let $u$ and $v$ be the eigenvectors of $A$ corresponding to the eigenvalues 1 and 3 respectively. Prove that $u+v$ is not an eigenvector of $A$.

\item  Let $A$ and $B$ be real matrices such that the sum of each row of $A$ is 1 and the sum of each row of $B$ is 2 . Then show that 2 is an eigenvalue of $A B$.
\end{enumerate}
 \hfill [JAM 2006]



\item  Suppose $W_{1}$ and $W_{2}$ are subspaces of $\mathbb{R}^{4}$ spanned by $\{(1,2,3,4),(2,1,1,2)\}$ and $\{(1,0,1,0),(3,0,1,0)\}$ respectively. Find a basis of $W_{1} \cap W_{2}$. Also find a basis of $W_{1}+W_{2}$ containing $\{(1,0,1,0),(3,0,1,0)\}$
 \hfill [JAM 2006]
 
  \item Determine $y_{0}$ such that the solution of the differential equation
$$
y^{\prime}-y=1-e^{-x}, y(0)=\mathrm{y}_{0}
$$
has a finite limit as $x \rightarrow \infty$.
 \hfill [JAM 2006]





  \item Let $\phi(x, y, z)=e^{x} \sin y$. Evaluate the surface integral $\iint_{S} \frac{\partial \phi}{\partial n} d \sigma$, where $S$ is the surface of the cube $0 \leq x \leq 1,0 \leq y \leq 1,0 \leq z \leq 1$ and $\frac{\partial \phi}{\partial n}$ is the directional derivative of $\phi$ in the direction of the unit outward normal to $S$. Verify the divergence theorem.
 \hfill [JAM 2006]

 
  \item Let $y=f(x)$ be a twice continuously differentiable function on $(0, \infty)$ satisfying
$$
f(1)=1 \text { and } f^{\prime}(x)=\frac{1}{2} f\left(\frac{1}{x}\right), x>0 .
$$
Form the second order differential equation satisfied by $y=f(x)$, and obtain its solution satisfying the given conditions.
 \hfill [JAM 2006]

  \item Let $G=\left\{\left(\begin{array}{ll}a & b \\ c & d\end{array}\right): a, b, c, d \in \mathbf{Z}\right\}$ be the group under matrix addition and $H$ be the subgroup of $G$ consisting of matrices with even entries. Find the order of the quotient group $G / H$.
   \hfill [JAM 2006]

  \item Let
$$
f(x)=\left\{\begin{array}{cc}
x^{2} & 0 \leq x \leq 1 \\
\sqrt{x} & x>1
\end{array}\right.
$$
Show that $f$ is uniformly continuous on $[0, \infty)$.




  \item Find $M_{n}=\max _{x \geq 0}\left\{\frac{x}{n\left(1+n x^{3}\right)}\right\}$, and hence prove that the series
$$
\sum_{n=1}^{\infty} \frac{x}{n\left(1+n x^{3}\right)}
$$
is uniformly convergent on $[0, \infty)$.


\item Let $R$ be the ring of polynomials with real coefficients under polynomial addition and polynomial multiplication. Suppose
$$
I=\{p \in R \text { : sum of the coefficients of } p \text { is zero }\} .
$$
Prove that $I$ is a maximal ideal of $R$.

\end{enumerate}
 \hfill [JAM 2006]





\newpage

\section{JAM 2007}


\begin{enumerate}
  \item Let $A(t)$ denote the area bounded by the curve $y=e^{-|x|}$, the $x$-axis and the straight lines $x=-t$ and $x=t$. Then $\lim _{t \rightarrow \infty} A(t)$ is equal to
  \begin{enumerate}
    \item 2
    \item 1
    \item $1 / 2$
    \item 0
\end{enumerate}
 \hfill [JAM 2007]


 
\item If $k$ is a constant such that $x y+k=e^{(x-1)^{2} / 2}$ satisfies the differential equation $x \frac{d y}{d x}=\left(x^{2}-x-1\right) y+(x-1)$, then $k$ is equal to
 \begin{enumerate}
\item 1
\item 0
\item $-1$
\item $-2$
\end{enumerate}
 \hfill [JAM 2007]



 
  \item Which of the following functions is uniformly continuous on the domain as stated?
  \begin{enumerate}
    \item $f(x)=x^{2}, \quad x \in \mathbf{R}$
    \item $f(x)=\frac{1}{x}, \quad x \in[1, \infty)$
    \item $f(x)=\tan x, \quad x \in(-\pi / 2, \pi / 2)$
    \item $f(x)=[x], \quad x \in[0,1]$
\end{enumerate}
 
( $[x]$ is the greatest integer less than or equal to $x$ )
\hfill [JAM 2007]




  \item Let $R$ be the ring of polynomials over $\mathbb{Z}_{2}$ and let $I$ be the ideal of $R$ generated by the polynomial $x^{3}+x+1$. Then the number of elements in the quotient ring $R / I$ is
    \begin{enumerate}
\item 2
\item 4
\item 8
\item 16
\end{enumerate}
 \hfill [JAM 2007]




\item Which of the following sets is a basis for the subspace
$W=\left\{\left[\begin{array}{ll}x & y \\ 0 & t\end{array}\right]: x+2 y+t=0, y+t=0\right\}$
of the vector space of all real $2 \times 2$ matrices?
\begin{enumerate}
    
\item $\left\{\left[\begin{array}{ll}1 & 0 \\ 0 & 0\end{array}\right],\left[\begin{array}{ll}0 & 1 \\ 0 & 0\end{array}\right],\left[\begin{array}{ll}0 & 0 \\ 0 & 1\end{array}\right]\right\}$

\item $\left\{\left[\begin{array}{cc}2 & 1 \\ 0 & -1\end{array}\right],\left[\begin{array}{cc}1 & -1 \\ 0 & 1\end{array}\right]\right\}$

\item $\left\{\left[\begin{array}{cc}-1 & 1 \\ 2 & -1\end{array}\right]\right\}$

\item $\left\{\left[\begin{array}{cc}1 & -1 \\ 0 & 1\end{array}\right]\right\}$
\end{enumerate}
 \hfill [JAM 2007]




\item Let $G$ be an Abelian group of order $10$. Let $S=\left\{g \in G: g^{-1}=g\right\}$. Then the number of non-identity elements in $S$ is
\begin{enumerate}
\item 5
\item 2
\item 1
\item 0
\end{enumerate}
 \hfill [JAM 2007]


 
\item  Let $\left(a_{n}\right)$ be an increasing sequence of positive real numbers such that the series $\sum_{k=1}^{\infty} a_{k}$ is divergent. Let $s_{n}=\sum_{k=1}^{n} a_{k}$ for $n=1,2, \ldots$ and $t_{n}=\sum_{k=2}^{n} \frac{a_{k}}{s_{k-1} s_{k}}$ for $n=2,3, \ldots$. Then $\lim _{n \rightarrow \infty} t_{n}$ is equal to
\begin{enumerate}
\item $1 / a_{1}$
\item 0
\item $1 /\left(a_{1}+a_{2}\right)$
\item $a_{1}+a_{2}$
\end{enumerate}
 \hfill [JAM 2007]




 
\item For every function $f:[0,1] \rightarrow \mathbf{R}$ which is twice differentiable and satisfies $f^{\prime}(x) \geq 1$ for all $x \in[0,1]$, we must have
\begin{enumerate}
  
\item $f^{\prime \prime}(x) \geq 0$ for all $x \in[0,1]$
\item $f(x) \geq x$ for all $x \in[0,1]$
\item $f\left(x_{2}\right)-x_{2} \leq f\left(x_{1}\right)-x_{1}$ for all $x_{1}, x_{2} \in[0,1]$ with $x_{2} \geq x_{1}$
\item $f\left(x_{2}\right)-x_{2} \geq f\left(x_{1}\right)-x_{1}$ for all $x_{1}, x_{2} \in[0,1]$ with $x_{2} \geq x_{1}$
\end{enumerate}
 \hfill [JAM 2007]


 

\item Let $f: \mathbf{R}^{2} \rightarrow \mathbf{R}$ be defined by
$$
f(x, y)=\left\{\begin{array}{cc}
\frac{x^{2} y}{x^{4}+y^{2}} & \text { if }(x, y) \neq(0,0) \\
0 & \text { if }(x, y)=(0,0)
\end{array}\right.
$$
Which of the following statements holds regarding the continuity and the existence of partial derivatives of $f$ at $(0,0)$ ?
\begin{enumerate}
  
\item Both partial derivatives of $f$ exist at $(0,0)$ and $f$ is continuous at $(0,0)$
\item Both partial derivatives of $f$ exist at $(0,0)$ and $f$ is NOT continuous at $(0,0)$
\item One partial derivative of $f$ does NOT exist at $(0,0)$ and $f$ is continuous at $(0,0)$
\item One partial derivative of $f$ does NOT exist at $(0,0)$ and $f$ is NOT continuous at $(0,0)$
\end{enumerate}
 \hfill [JAM 2007]



 
\item Suppose $\left(c_{n}\right)$ is a sequence of real numbers such that $\lim _{n \rightarrow \infty}\left|c_{n}\right|^{1 / n}$ exists and is non-zero. If the radius of convergence of the power series $\sum_{n=0}^{\infty} c_{n} x^{n}$ is equal to $r$, then the radius of convergence of the power series $\sum_{n=1}^{\infty} n^{2} c_{n} x^{n}$ is

\begin{enumerate}
\item less than $r$
\item greater than $r$
\item equal to $r$
\item equal to 0
\end{enumerate}
 \hfill [JAM 2007]


 
\item The rank of the matrix $\left[\begin{array}{ccc}1 & 4 & 8 \\ 2 & 10 & 22 \\ 0 & 4 & 12\end{array}\right]$ is
\begin{enumerate}
\item 3
\item 2
\item 1
\item 0
\end{enumerate}
 \hfill [JAM 2007]





  \item Let $f: \mathbf{R} \rightarrow \mathbf{R}$ be a continuous function. If $\int_{0}^{x} f(2 t) d t=\frac{x}{\pi} \sin (\pi x)$ for all $x \in \mathbf{R}$, then $f(2)$ is equal to
  \begin{enumerate}
\item $-1$
\item 0
\item 1
\item 2
\end{enumerate}
 \hfill [JAM 2007]



 
\item Let $\vec{u}=\left(a e^{x} \sin y-4 x\right) \hat{\mathrm{i}}+\left(2 y+e^{x} \cos y\right) \hat{\mathrm{j}}+\alpha z \hat{\mathrm{k}}$, where $a$ is a constant. If the line integral $\oint_{C} \vec{u} \cdot d \vec{r}$ over every closed curve $C$ is zero, then $a$ is equal to
\begin{enumerate}
\item $-2$
\item $-1$
\item $0$
\item $1$
\end{enumerate}
 \hfill [JAM 2007]

\item One of the integrating factors of the differential equation $\left(y^{2}-3 x y\right) d x+\left(x^{2}-x y\right) d y=0$ is
\begin{enumerate}
\item $1 /\left(x^{2} y^{2}\right)$
\item $1 /\left(x^{2} y\right)$
\item $1 /\left(x y^{2}\right)$
\item $1 /(x y)$
\end{enumerate}
 \hfill [JAM 2007]



  \item Let $C$ denote the boundary of the semi-circular disk $D=\left\{(x, y) \in \mathbf{R}^{2}: x^{2}+y^{2} \leq 1, y \geq 0\right\}$ and let $\varphi(x, y)=x^{2}+y$ for $(x, y) \in D$. If $\hat{n}$ is the outward unit normal to $C$, then the integral $\oint_{C}(\vec{\nabla} \varphi) \cdot \hat{n} d s$, evaluated counter-clockwise over $C$, is equal to
  \begin{enumerate}
\item 0
\item $\pi-2$
\item $\pi$
\item $\pi+2$ 
\end{enumerate}
 \hfill [JAM 2007]

 
\item 
\begin{enumerate}
    \item Let $M=\left[\begin{array}{ccc}1+i & 2 i & i+3 \\ 0 & 1-i & 3 i \\ 0 & 0 & i\end{array}\right]$. Determine the eigenvalues of the matrix $B=M^{2}-2 M+I$

\item  Let $N$ be a square matrix of order 2 . If the determinant of $N$ is equal to 9 and the sum of the diagonal entries of $N$ is equal to $10$ , then determine the eigenvalues of $N$.
\end{enumerate} 
 \hfill [JAM 2007]





\item 
\begin{enumerate}
  \item Using the method of variation of parameters, solve the differential equation

$$
x^{2} \frac{d^{2} y}{d x^{2}}+x \frac{d y}{d x}-y=x^{2},
$$
given that $x$ and $\frac{1}{x}$ are two solutions of the corresponding homogeneous equation.

\item  Find the real number $\alpha$ such that the differential equation
$$
\frac{d^{2} y}{d x^{2}}+2(\alpha-1)(\alpha-3) \frac{d y}{d x}+(\alpha-2) y=0
$$
has a solution $y(x)=a \cos (\beta x)+b \sin (\beta x)$ for some non-zero real numbers $a, b, \beta$.
\end{enumerate} 
 \hfill [JAM 2007]



 \item 
\begin{enumerate}
  \item Let $a, b, c$ be non-zero real numbers such that $(a-b)^{2}=4 a c$. Solve the differential equation $a(x+\sqrt{2})^{2} \frac{d^{2} y}{d x^{2}}+b(x+\sqrt{2}) \frac{d y}{d x}+c y=0$.

\item Solve the differential equation
$$
d x+\left(e^{y \sin y}-x\right)(y \cos y+\sin y) d y=0 .
$$
\end{enumerate}
 \hfill [JAM 2007]




  \item Let $f(x, y)=x\left(x-2 y^{2}\right)$ for $(x, y) \in \mathbf{R}^{2}$. Show that $f$ has a local minimum at $(0,0)$ on every straight line through $(0,0)$. Is $(0,0)$ a critical point of $f$ ? Find the discriminant of $f$ at $(0,0)$. Does $f$ have a local minimum at $(0,0)$ ? Justify your answers.
 \hfill [JAM 2007]
  
  
  
  
\item 
\begin{enumerate}
  \item  Find the finite volume enclosed by the paraboloids $z=2-x^{2}-y^{2}$ and $z=x^{2}+y^{2}$.
  \item  Let $f:[0,3] \rightarrow \mathbf{R}$ be a continuous function with $\int_{0}^{3} f(x) d x=3$. Evaluate
$$
\int_{0}^{3}\left[x f(x)+\int_{0}^{x} f(t) d t\right] d x
$$
\end{enumerate}
 \hfill [JAM 2007]




\item 
\begin{enumerate}
  \item  Let $S$ be the surface $\left\{(x, y, z) \in \mathbf{R}^{3}: x^{2}+y^{2}+2 z=2, z \geq 0\right\}$, and let $\hat{n}$ be the outward unit normal to $S$. If $\vec{F}=y \hat{\mathrm{i}}+x z \hat{\mathrm{j}}+\left(x^{2}+y^{2}\right) \hat{\mathrm{k}}$, then evaluate the integral $\iint_{S} \vec{F} \cdot \hat{n} d S$.

\item  Let $\vec{r}=x \hat{i}+y \hat{j}+z \hat{k}$ and $r=|\vec{r}|$. If a scalar field $\varphi$ and a vector field $\vec{u}$ satisfy $\vec{\nabla} \varphi=\vec{\nabla} \times \vec{u}+f(r) \vec{r}$, where $f$ is an arbitrary differentiable function, then show that $\nabla^{2} \varphi=r f^{\prime}(r)+3 f(r)$
\end{enumerate}
 \hfill [JAM 2007]


\item 
\begin{enumerate}
  \item Let $D$ be the region bounded by the concentric spheres $S_{1}: x^{2}+y^{2}+z^{2}=a^{2}$ and $S_{2}: x^{2}+y^{2}+z^{2}=b^{2}$, where $a<b$. Let $\hat{n}$ be the unit normal to $S_{1}$ directed away from the origin. If $\nabla^{2} \varphi=0$ in $D$ and $\varphi=0$ on $S_{2}$, then show that $\iiint_{D}|\vec{\nabla} \varphi|^{2} d V+\iint_{S_{1}} \varphi(\vec{\nabla} \varphi) \cdot \hat{n} d S=0$

\item  Let $C$ be the curve in $\mathbf{R}^{3}$ given by $x^{2}+y^{2}=a^{2}, z=0$ traced counter-clockwise, and let $\vec{F}=x^{2} y^{3} \hat{\mathrm{i}}+\hat{\mathrm{j}}+z \hat{\mathrm{k}}$. Using Stokes' theorem, evaluate $\oint_{C} \vec{F} \cdot d \vec{r}$.
\end{enumerate}
 \hfill [JAM 2007]





  \item Let $V$ be the subspace of $\mathbf{R}^{4}$ spanned by the vectors $(1,0,1,2),(2,1,3,4)$ and $(3,1,4,6)$. Let $T: V \rightarrow \mathbf{R}^{2}$ be a linear transformation given by $T(x, y, z, t)=(x-y, z-t)$ for all $(x, y, z, t) \in V$. Find a basis for the null space of $T$ and also a basis for the range space of $T$
 \hfill [JAM 2007]


\item   
\begin{enumerate}
    \item  Compute the double integral $\iint_{D}(x+2 y) d x d y$, where $D$ is the region in the $x y$-plane bounded by the straight lines $y=x+3, y=x-3, y=-2 x+4$ and $y=-2 x-2$.

    \item  Evaluate $\int_{0}^{\pi / 2}\left[\int_{\pi / 2}^{\pi} \frac{\sin x}{x} d x\right] d y+\int_{\pi / 2}^{\pi}\left[\int_{y}^{\pi} \frac{\sin x}{x} d x\right] d y$.
\end{enumerate}
 \hfill [JAM 2007]


\item 
\begin{enumerate}
    \item  Does the series $\sum_{k=1}^{\infty} \frac{(-1)^{k} k+x^{k}}{k^{2}}$ converge uniformly for $x \in[-1,1]$ ? Justify.


    \item  Suppose $\left(f_{n}\right)$ is a sequence of real-valued functions defined on $\mathbf{R}$ and $f$ is a real-valued function defined on $\mathbf{R}$ such that $\left|f_{n}(x)-f(x)\right| \leq\left|a_{n}\right|$ for all $n \in \mathbf{N}$ and $a_{n} \rightarrow 0$ as $n \rightarrow \infty$. Must the sequence $\left(f_{n}\right)$ be uniformly convergent on $\mathbf{R}$ ? Justify. (6) 
\end{enumerate}
\hfill [JAM 2007] 
    
\item 
\begin{enumerate}
    \item  Suppose $f$ is a real-valued thrice differentiable function defined on $\mathbf{R}$ such that $f^{\prime \prime \prime}(x)>0$ for all $x \in \mathbf{R}$. Using Taylor's formula, show that $f\left(x_{2}\right)-f\left(x_{1}\right)>\left(x_{2}-x_{1}\right) f^{\prime}\left(\frac{x_{1}+x_{2}}{2}\right)$ for all $x_{1}$ and $x_{2}$ in $\mathbf{R}$ with $x_{2}>x_{1}$.

\item  Let $\left(a_{n}\right)$ and $\left(b_{n}\right)$ be sequences of real numbers such that $a_{n} \leq a_{n+1} \leq b_{n+1} \leq b_{n}$ for all $n \in \mathbf{N}$. Must there exist a real number $x$ such that $a_{n} \leq x \leq b_{n}$ for all $n \in \mathbf{N}$ ? Justify your answer.
\end{enumerate}
\hfill [JAM 2007] 



  \item Let $G$ be the group of all $2 \times 2$ matrices with real entries with respect to matrix multiplication. Let $G_{1}$ be the smallest subgroup of $G$ containing $\left[\begin{array}{ll}0 & 1 \\ 1 & 0\end{array}\right]$ and $\left[\begin{array}{cc}-1 & 0 \\ 0 & -1\end{array}\right]$, and $G_{2}$ be the smallest subgroup of $G$ containing $\left[\begin{array}{cc}0 & -1 \\ 1 & 0\end{array}\right]$ and $\left[\begin{array}{cc}0 & 1 \\ -1 & 0\end{array}\right]$. Determine all elements of $G_{1}$ and find their orders. Determine all elements of $G_{2}$ and find their orders. Does there exist a one-to-one homomorphism from $G_{1}$ onto $G_{2}$ ? Justify.
 \hfill [JAM 2007] 

  
\item 
\begin{enumerate}
    \item Let $p$ be a prime number and let $\mathbf{Z}$ be the ring of integers. If an ideal $J$ of $\mathbf{Z}$ contains the set $p \mathbf{Z}$ properly, then show that $J=\mathbf{Z}$. (Here $p \mathbf{Z}=\{p x: x \in \mathbf{Z}\}$.)


    \item  Consider the ring $R=\{a+i b: a, b \in \mathbf{Z}\}$ with usual addition and multiplication. Find all invertible elements of $R$.
\end{enumerate}
\hfill [JAM 2007] 


\item 
\begin{enumerate}
    \item  Suppose $E$ is a non-empty subset of $\mathbf{R}$ which is bounded above, and let $\alpha=\sup E$. If $E$ is closed, then show that $\alpha \in E$. If $E$ is open, then show that $\alpha \notin E$.

    \item  Find all limit points of the set $E=\left\{n+\frac{1}{2 m}: n, m \in \mathbf{N}\right\}$.
\end{enumerate}
\hfill [JAM 2007] 


\end{enumerate}




\newpage
\section{JAM 2008}

\begin{enumerate}
   
\item  The least positive integer $n$, such that $\left(\begin{array}{cc}\cos \frac{\pi}{4} & \sin \frac{\pi}{4} \\ -\sin \frac{\pi}{4} & \cos \frac{\pi}{4}\end{array}\right)^{n}$ is the identity matrix of order 2, is
\begin{enumerate}
\item 4  
\item 8 
\item 12 
\item 16
\end{enumerate}
 \hfill [JAM 2008]


\item  Let $S=\left\{T: \mathbb{R}^{3} \rightarrow \mathbb{R}^{3} \mid T\right.$ is a linear transformation with $T(1,0,1)=(1,2,3)$, $T(1,2,3)=(1,0,1)\}$. Then $S$ is 
\begin{enumerate}
\item a singleton set 
\item a countably infinite set 
\item a finite set containing more than one element 
\item an uncountable set
\end{enumerate}
 \hfill [JAM 2008]


\item  Let $s_{n}=\int_{0}^{1} \frac{n x^{n-1}}{(1+x)} d x$ for $n \geq 1$. Then as $n \rightarrow \infty$, the sequence $\left\{s_{n}\right\}$ tends to 

\begin{enumerate}
\item 0 
\item $1 / 2$ 
\item 1 
\item $+\infty$
\end{enumerate}
 \hfill [JAM 2008]


\item  The work done by the force $\vec{F}=4 y \hat{i}-3 x y \hat{j}+z^{2} \hat{k}$ in moving a particle over the circular path $x^{2}+y^{2}=1, z=0$ from $(1,0,0)$ to $(0,1,0)$ is 
\begin{enumerate}
\item $\pi+1$ 
\item $\pi-1$ 
\item $-\pi+1$ 
\item $-\pi-1$
\end{enumerate}
 \hfill [JAM 2008]


\item  The set of all boundary points of $\mathbb{Q}$ in $\mathbb{R}$ is 
\begin{enumerate}
\item $\mathbb{R}$ 
\item $\mathbb{R} \backslash \mathbb{Q}$ 
\item $\mathbb{Q}$ 
\item $\varnothing$
\end{enumerate}
 \hfill [JAM 2008]


\item  Let $V=\left\{(x, y, z) \in \mathbb{R}^{3}: \frac{1}{4} \leq x^{2}+y^{2}+z^{2} \leq 1\right\}$ and $\vec{F}=\frac{x \hat{i}+y \hat{j}+z \hat{k}}{\left(x^{2}+y^{2}+z^{2}\right)^{2}}$ for $(x, y, z) \in V$. Let $\hat{n}$ denote the outward unit normal vector to the boundary of $V$ and $S$ denote the part $\left\{(x, y, z) \in \mathbb{R}^{3}: x^{2}+y^{2}+z^{2}=\frac{1}{4}\right\}$ of the boundary of $V$. Then $\iint_{S} \vec{F} \cdot \hat{n} d S$ is equal to 
\begin{enumerate}
\item $-8 \pi$ 
\item $-4 \pi$ 
\item $4 \pi$ 
\item $8 \pi$
\end{enumerate}
 \hfill [JAM 2008]


\item  The set $U=\left\{x \in \mathbb{R} \mid \sin x=\frac{1}{2}\right\}$ is 
\begin{enumerate}
\item open 
\item both open and closed 
\item closed 
\item neither open nor closed
\end{enumerate}
 \hfill [JAM 2008]


\item  Let $f(x)=\int_{0}^{x}\left(x^{2}-t^{2}\right) g(t) d t$, where $g$ is a real valued continuous function on $\mathbb{R}$. Then $f^{\prime}(x)$ is equal to 
\begin{enumerate}
\item 0 
\item $x^{3} g(x)$ 
\item $\int_{0}^{x} g(t) d t$ 
\item $2 x \int_{0}^{x} g(t) d t$
\end{enumerate}
 \hfill [JAM 2008]




 \item  Let $y_{1}(x)$ and $y_{2}(x)$ be linearly independent solutions of the differential equation $y^{\prime \prime}+P(x) y^{\prime}+Q(x) y=0$, where $P(x)$ and $Q(x)$ are continuous functions on an interval $I$. Then $y_{3}(x)=a y_{1}(x)+b y_{2}(x)$ and $y_{4}(x)=c y_{1}(x)+d y_{2}(x)$ are linearly independent solutions of the given differential equation if 
 \begin{enumerate}
\item $a d=b c$ 
\item $a c=b d$ 
\item $a d \neq b c$ 
\item $a c \neq b d$

\end{enumerate}
 \hfill [JAM 2008]

\item  The set $R=\{f \mid f$ is a function from $\mathbb{Z}$ to $\mathbb{R}\}$ under the binary operations $+$ and . defined as $(f+g)(n)=f(n)+g(n)$ and $(f \cdot g)(n)=f(n) g(n)$ for all $n \in \mathbb{Z}$ forms a ring. Let $S_{1}=\{f \in R \mid f(-n)=f(n)$ for all $n \in \mathbb{Z}\}$ and $S_{2}=\{f \in R \mid f(0)=0\}$. Then 
\begin{enumerate}
\item $S_{1}$ and $S_{2}$ are both ideals in $R$ 
\item $S_{2}$ is an ideal in $R$ while $S_{1}$ is not 
\item $S_{1}$ is an ideal in $R$ while $S_{2}$ is not 
\item neither $S_{1}$ nor $S_{2}$ is an ideal in $R$
\end{enumerate}
 \hfill [JAM 2008]




\item  Let $T: \mathbb{R}^{3} \rightarrow \mathbb{R}^{3}$ be a linear transformation such that $T(1,2,3)=(1,2,3), T(1,5,0)=(2,10,0)$ and $T(-1,2,-1)=(-3,6,-3)$. The dimension of the vector space spanned by all the eigenvectors of $T$ is 
\begin{enumerate}
\item 0 
\item 1 
\item 2 
\item 3
\end{enumerate}
 \hfill [JAM 2008]


\item  Let $\left\{a_{n}\right\}$ and $\left\{b_{n}\right\}$ be sequences of real numbers defined as $a_{1}=1$ and for $n \geq 1$, $a_{n+1}=a_{n}+(-1)^{n} 2^{-n}, \quad b_{n}=\frac{2 a_{n+1}-a_{n}}{a_{n}}$. Then 
\begin{enumerate}
\item $\left\{a_{n}\right\}$ converges to zero and $\left\{b_{n}\right\}$ is a Cauchy sequence 
\item $\left\{a_{n}\right\}$ converges to a non-zero number and $\left\{b_{n}\right\}$ is a Cauchy sequence 
\item $\left\{a_{n}\right\}$ converges to zero and $\left\{b_{n}\right\}$ is not a convergent sequence 
\item $\left\{a_{n}\right\}$ converges to a non-zero number and $\left\{b_{n}\right\}$ is not a convergent sequence
\end{enumerate}
 \hfill [JAM 2008]


\item  Let $f:(-1,1) \rightarrow \mathbb{R}$ be defined as $f(x)=\frac{x^{2}}{1-\cos x}$ for $x \neq 0$ and $f(0)=2$. If $f(x)=\sum_{n=0}^{\infty} a_{n} x^{n}$ is the Taylor expansion of $f$ for all $x$ in $(-1,1)$, then $\sum_{n=0}^{x} a_{2 n+1}$ is 
\begin{enumerate}
\item 0 
\item $1 / 2$ 
\item 1 
\item 2
\end{enumerate}
 \hfill [JAM 2008]


\item  Let $y_{1}(x)$ and $y_{2}(x)$ be twice differentiable functions on an interval $I$ satisfying the differential equations $\frac{d y_{1}}{d x}-y_{1}-y_{2}=e^{x}$ and $2 \frac{d y_{1}}{d x}+\frac{d y_{2}}{d x}-6 y_{1}=0$. Then $y_{1}(x)$ is 
\begin{enumerate}
\item $C_{1} e^{-2 x}+C_{2} e^{3 x}-\frac{1}{4} e^{x}$ 
\item $C_{1} e^{2 x}+C_{2} e^{-3 x}+\frac{1}{4} e^{x}$ 
\item $C_{1} e^{2 x}+C_{2} e^{-3 x}-\frac{1}{4} e^{x}$ 
\item $C_{1} e^{-2 x}+C_{2} e^{3 x}+\frac{1}{4} e^{x}$
\end{enumerate}
 \hfill [JAM 2008]


\item  Let $G$ be a finite group and $H$ be a normal subgroup of $G$ of order 2. Then the order of the center of $G$ is 
\begin{enumerate}
\item 0 
\item 1 
\item an even integer $\geq 2$ 
\item an odd integer $\geq 3$
\end{enumerate}
 \hfill [JAM 2008]




\item 
\begin{enumerate}
\item Let $f$ and $g$ be continuous functions on $\mathbb{R}$ such that $f(x)=\int_{0}^{x} g(t) d t$ and
$$
g(x)=\int_{x}^{0} f(t) d t+1 \text {. Prove that }(f(x))^{2}+(g(x))^{2}=1 \text { for all } x \in \mathbb{R}
$$
\item  Let $f: \mathbb{R} \rightarrow \mathbb{R}$ be a function such that $f^{\prime}$ is continuous on $\mathbb{R}$. Show that the series
$$
\sum_{n=1}^{x}\left(f\left(\frac{x}{2 n}\right)-f\left(\frac{x}{2 n+1}\right)\right) \text { converges uniformly on }[0,1] \text {. }
$$
\end{enumerate}
 \hfill [JAM 2008]


\item 
\begin{enumerate}
\item   Find the maxima, minima and saddle points, if any, for the function $f(x, y)=\left(y-x^{2}\right)\left(y-2 x^{2}\right)$ on $\mathbb{R}^{2}$

\item  Let $P(x)=a_{0}+a_{1} x^{2}+a_{2} x^{4}+\cdots+a_{n} x^{2 n}$, where $n \geq 1$ and $a_{k}>0$ for $k=0$, $1, \cdots, n$. Show that $P(x)-x P^{\prime}(x)=0$ has exactly two real roots.
\end{enumerate}
 \hfill [JAM 2008]


\item 
\begin{enumerate}
\item   Given that $y_{1}(x)=x$ is a solution of $\left(1+x^{2}\right) y^{\prime \prime}-2 x y^{\prime}+2 y=0, x>0$, find a second linearly independent solution.

\item Solve $x^{2} y^{\prime \prime}+x y^{\prime}-y=4 x \log x, x>0$.

\end{enumerate}
 \hfill [JAM 2008]

\item 
\begin{enumerate}
\item   Let $\phi$ be a differentiable function on $[0,1]$ satisfying $\phi^{\prime}(x) \leq 1+3 \phi(x)$ for all $x \in[0,1]$ with $\phi(0)=0$. Show that $3 \phi(x) \leq e^{3 x}-1$.

\item  If $y_{1}(x)=x(1-2 x), y_{2}(x)=2 x(1-x)$ and $y_{3}(x)=x\left(e^{x}-2 x\right)$ are three solutions of a non-homogeneous linear differential equation $y^{\prime \prime}+P(x) y^{\prime}+Q(x) y=R(x)$, where $P(x), Q(x)$ and $R(x)$ are continuous functions on $[a, b]$ with $a>0$, then find its general solution.
\end{enumerate}
 \hfill [JAM 2008]


\item 
\begin{enumerate}
\item  
$$
\text { Evaluate } \int_{1}^{4} \int_{0}^{1} \int_{2 y}^{2} \frac{\cos x^{2}}{\sqrt{z}} d x d y d z
$$
\item  Find the surface area of the portion of the cone $z^{2}=x^{2}+y^{2}$ that is inside the cylinder $z^{2}=2 y$
\end{enumerate}
 \hfill [JAM 2008]


\item 

\begin{enumerate}
\item   Use Green's theorem to evaluate the integral $\oint_{C} x^{2} d x+\left(x+y^{2}\right) d y$, where $C$ is the closed curve given by $y=0, y=x$ and $y^{2}=2-x$ in the first quadrant, oriented counter clockwise.

\item  Let $f: \mathbb{R} \rightarrow \mathbb{R}$ be a continuous function. Use change of variables to prove that
$$
\iint_{D} f(x-y) d x d y=\int_{-1}^{1} f(u) d u \text {, where } D=\left\{(x, y) \in \mathbb{R}^{2}:|x|+|y| \leq 1\right\} \text {. }
$$
\end{enumerate}
 \hfill [JAM 2008]




\item  Using Gauss's divergence theorem, evaluate the integral $\iint_{S} \vec{F} \cdot \hat{n} d S$, where $\vec{F}=4 x z \hat{i}-y^{2} \hat{j}+4 y z \hat{k}, \quad S$ is the surface of the solid bounded by the sphere $x^{2}+y^{2}+z^{2}=10$ and the paraboloid $x^{2}+y^{2}=z-2$, and $\hat{n}$ is the outward unit normal vector to $S$.
\hfill [JAM 2008]



\item
\begin{enumerate}
 \item A square matrix $M$ of order $n$ with complex entries is called skew Hermitian if $M+\bar{M}^{T}=0$, where 0 is the zero matrix of order $n$.

Determine whether $V=\{M \mid M$ is a $2 \times 2$ skew Hermitian matrix $\}$ is a vector space over (i) the field $\mathbb{R}$ and (ii) the field $\mathbb{C}$ with the usual operations of addition and scalar multiplication for matrices?

\item  Let $V=\{P(x) \mid P(x)$ is a polynomial of degree $\leq n$ with real coefficients $\}$ and $T: V \rightarrow \mathbb{R}^{m}$ be defined as $T(P(x))=(P(1), P(2), \cdots, P(m))$. Then show that $T$ is linear and determine the Nullity of $T$.
\end{enumerate}
 \hfill [JAM 2008]


\item  Let $G$ be the set of all $3 \times 3$ real matrices $M$ such that $M M^{T}=M^{T} M=I_{3}$ and let $H=\{M \in G \mid \operatorname{det} M=1\}$, where $I_{3}$ is the identity matrix of order 3 . Then show that
\begin{enumerate}
\item  $G$ is a group under matrix multiplication,

\item  $H$ is a normal subgroup of $G$,

\item  $\phi: G \rightarrow\{-1,1\}$ given by $\phi(M)=\operatorname{det} M$ is onto,

\item $G / H$ is abelian.
\end{enumerate}
 \hfill [JAM 2008]


\item 
\begin{enumerate}
\item Suppose that $(R,+, \cdot)$ is a ring having the property $a \cdot b=c \cdot a \Rightarrow b=c$, when $a \neq 0$. Then prove that $(R,+, \cdot)$ is a commutative ring.

\item  Let $R$ be a commutative ring with identity. For $a_{1}, a_{2}, \cdots, a_{n} \in R$, the ideal generated by $\left\{a_{1}, a_{2}, \cdots, a_{n}\right\}$ is given by
$$
\left\langle a_{1}, a_{2}, \cdots, a_{n}\right\rangle=\left\{r_{1} a_{1}+r_{2} a_{2}+\cdots+r_{n} a_{n} \mid r_{i} \in R, 1 \leq i \leq n\right\} .
$$
Let $\mathbb{Z}[x]$ be the set of all polynomials with integer coefficients. Consider the ideal $I=\{f \in \mathbb{Z}[x] \mid f(0)$ is an even integer $\}$. Prove that $I=\langle 2, x\rangle$ and that it is a maximal ideal.
\end{enumerate}
 \hfill [JAM 2008]




\item  For a given positive integer $n>1$, show that there exist subspaces $X_{1}, X_{2}, \cdots, X_{n}$ of $\mathbb{R}^{m}$ for some integer $m>n$ and a linear transformation $T: \mathbb{R}^{m} \rightarrow \mathbb{R}^{m}$ such that

\begin{itemize}
  \item $\operatorname{dim} X_{k}=k, k=1,2, \cdots, n$,

  \item for $i \neq j, X_{i} \cap X_{j}=\{0\}$, where $\overrightarrow{0}$ is the zero vector of $\mathbb{R}^{m}$,

  \item $T\left(X_{k}\right)=X_{k-1}, k=1,2, \cdots, n$, where $X_{0}=\{\overrightarrow{0}\}$.

\end{itemize}
Also, find the rank of $T$.
\hfill [JAM 2008]


\item  Let $f:(0, \infty) \rightarrow(0, \infty)$ be a continuously differentiable function and let $z=\frac{x y}{f\left(x^{2}+y^{2}\right)}$ be defined for $x y \neq 0$.

\begin{enumerate}
\item Prove that $\frac{\partial z}{\partial x}+\frac{\partial z}{\partial y}=\frac{(x+y)}{\left[f\left(x^{2}+y^{2}\right)\right]^{2}}\left\{f\left(x^{2}+y^{2}\right)-2 x y f^{\prime}\left(x^{2}+y^{2}\right)\right\}$.

\item  Further, if $f$ is homogeneous of degree $\frac{1}{2}$, then verify that $x \frac{\partial z}{\partial x}+y \frac{\partial z}{\partial y}=z$.
\end{enumerate}
 \hfill [JAM 2008]
 

\item  Determine the interval of convergence of the power series $\sum_{n=1}^{x} n(2 n-1) x^{2 n}$ and show that its sum is $\frac{x^{2}\left(1+3 x^{2}\right)}{\left(1-x^{2}\right)^{3}}$ at any point $x$ in its interval of convergence.
\hfill [JAM 2008]



\item
\begin{enumerate}
\item Let $f: \mathbb{R}^{2} \rightarrow \mathbb{R}$ be defined as $f(x, y)=x^{2} \cos (y / x)$ for $x \neq 0$ and $f(x, y)=0$ for $x=0$. Compute $\frac{\partial f}{\partial x}$ at all points in $\mathbb{R}^{2}$ and show that it is continuous at the origin.

\item  Let $f:(0,1) \rightarrow(0, \infty)$ be a uniformly continuous function. If $\left\{x_{n}\right\}$ is a Cauchy sequence in $(0,1)$, then prove that $\left\{f\left(x_{n}\right)\right\}$ is a Cauchy sequence in $(0, \infty)$. Hence deduce that for any two Cauchy sequences $\left\{x_{n}\right\}$ and $\left\{y_{n}\right\}$ in $(0,1)$, $\left\{\left|f\left(x_{n}\right)-f\left(y_{n}\right)\right|\right\}$ is a Cauchy sequence in $(0, \infty)$.
\end{enumerate}
 \hfill [JAM 2008]





\end{enumerate}
\newpage
\section{JAM 2009}


\begin{enumerate}

\item  Let $V$ be the vector space of all $6 \times 6$ real matrices over the field $\mathrm{R}$. Then the dimension of the subspace of $V$ consisting of all symmetric matrices is 
\begin{enumerate}
   \item 15 
\item 18 
\item 21 
\item 35 
\end{enumerate}
\hfill [JAM 2009]




\item  Let $R$ be the ring of all functions from $R$ to $R$ under point-wise addition and multiplication. Let $I=\{f: \mathbb{R} \rightarrow \mathbb{R} \mid f$ is a bounded function $\}, J=\{f: \mathbb{R} \rightarrow \mathbb{R} \mid f(3)=0\}$. Then 
\begin{enumerate}
   
\item $J$ is an ideal of $R$ but $I$ is not an ideal of $R$ 
\item $I$ is an ideal of $R$ but $J$ is not an ideal of $R$ 
\item both $I$ and $J$ are ideals of $R$ 
\item neither $I$ nor $J$ is an ideal of $R$ 
\end{enumerate}
\hfill [JAM 2009]



\item  Which of the following sequences of functions is uniformly convergent on $(0,1)$ ? 
\begin{enumerate}
    \item $x^{n}$ 
\item $\frac{n}{n x+1}$ 
\item $\frac{x}{n x+1}$ 
\item $\frac{1}{n x+1}$
\end{enumerate}
\hfill [JAM 2009]




\item  Let $T: \mathrm{R}^{4} \rightarrow \mathrm{R}^{4}$ be a linear transformation satisfying $T^{3}+3 T^{2}=4 I$, where $I$ is the identity transformation. Then the linear transformation $S=T^{4}+3 T^{3}-4 I$ is 
\begin{enumerate}
    
\item one-one but not onto 
\item invertible 
\item onto but not one-one 
\item non-invertible
\end{enumerate}
\hfill [JAM 2009]



\item  The number of all subgroups of the group $\left(Z_{60},+\right)$ of integers modulo 60 is 
\begin{enumerate}
    \item 2 
\item 10 
\item 12 
\item 60
\end{enumerate}
\hfill [JAM 2009]




\item  Let $a_{n}= \begin{cases}\frac{1}{3^{n}} & \text { if } n \text { is a prime, } \\ \frac{1}{4^{n}} & \text { if } n \text { is not aprime. }\end{cases}$

Then the radius of convergence of the power series $\sum_{n=1}^{\infty} a_{n} x^{n}$ is 
\begin{enumerate}
    \item 4 
\item 3 
\item $\frac{1}{3}$ 
\item $\frac{1}{4}$
\end{enumerate}
\hfill [JAM 2009]




\item  The set of all limit points of the sequence $1, \frac{1}{2}, \frac{1}{4}, \frac{3}{4}, \frac{1}{8}, \frac{3}{8}, \frac{5}{8}, \frac{7}{8}, \frac{1}{16}, \frac{3}{16}, \frac{5}{16}, \frac{7}{16}, \frac{9}{16}, \ldots$ is 
\begin{enumerate}
    \item $[0,1]$ 
\item $(0,1]$ 
\item the set of all rational numbers in $[0,1]$ 
\item the set of all rational numbers in $[0,1]$ of the form $\frac{m}{2^{n}}$ where $m$ and $n$ are integers 
\end{enumerate}
\hfill [JAM 2009]


\item  Let $F: \mathrm{R} \rightarrow \mathbb{R}$ be a continuous function and $a>0$. Then the integral $\int_{0}^{a}\left[\int_{0}^{x} F(y) d y\right] d x$ equals
\begin{enumerate}
    \item $\int_{0}^{a} y F(y) d y$ 
\item $\int_{0}^{a}(a-y) F(y) d y$ 
\item $\int_{0}^{a}(y-a) F(y) d y$ 
\item $\int_{a}^{0} y F(y) d y$
\end{enumerate}
\hfill [JAM 2009]

 


\item  The set of all positive values of $a$ for which the series $\sum_{n=1}^{\infty}\left(\frac{1}{n}-\tan ^{-1}\left(\frac{1}{n}\right)\right)^{a}$ converges, is 
\begin{enumerate}
    
\item $\left(0, \frac{1}{3}\right]$ 
\item $\left(0, \frac{1}{3}\right)$ 
\item $\left[\frac{1}{3}, \infty\right)$ 
\item $\left(\frac{1}{3}, \infty\right)$
\end{enumerate}
\hfill [JAM 2009]



\item  Let $a$ be a non-zero real number. Then $\lim _{x \rightarrow a} \frac{1}{x^{2}-a^{2}} \int_{a}^{x} \sin \left(t^{2}\right) d t$ equals 
\begin{enumerate}
    
\item $\frac{1}{2 a} \sin \left(a^{2}\right)$ 
\item $\frac{1}{2 a} \cos \left(a^{2}\right)$ 
\item $-\frac{1}{2 a} \sin \left(a^{2}\right)$ 
\item $-\frac{1}{2 a} \cos \left(a^{2}\right)$
\end{enumerate}
\hfill [JAM 2009]



\item  Let $T(x, y, z)=x y^{2}+2 z-x^{2} z^{2}$ be the temperature at the point $(x, y, z)$. The unit vector in the direction in which the temperature decreases most rapidly at $(1,0,-1)$ is
\begin{enumerate}
    \item $-\frac{1}{\sqrt{5}} \hat{i}+\frac{2}{\sqrt{5}} \hat{k}$ 
\item $\frac{1}{\sqrt{5}} \hat{i}-\frac{2}{\sqrt{5}} \hat{k}$ 
\item $\frac{2}{\sqrt{14}} \hat{i}+\frac{3}{\sqrt{14}} \hat{j}+\frac{1}{\sqrt{14}} \hat{k}$ 
\item $-\left(\frac{2}{\sqrt{14}} \hat{i}+\frac{3}{\sqrt{14}} \hat{j}+\frac{1}{\sqrt{14}} \hat{k}\right)$

\end{enumerate}
\hfill [JAM 2009]

 

\item  Consider the differential equation $2 \cos \left(y^{2}\right) d x-x y \sin \left(y^{2}\right) d y=0$. Then 
\begin{enumerate}
    \item $e^{x}$ is an integrating factor 
\item $e^{-x}$ is an integrating factor 
\item $3 x$ is an integrating factor 
\item $x^{3}$ is an integrating factor
\end{enumerate}
\hfill [JAM 2009]




\item  Suppose $\vec{V}=p(x, y) \hat{i}+q(x, y) \hat{j}$ is a continuously differentiable vector field defined in a domain $D$ in $R^{2}$. Which one of the following statements is NOT equivalent to the remaining ones? 
\begin{enumerate}
    \item There exists a function $\phi(x, y)$ such that $\frac{\partial \phi}{\partial x}=p(x, y)$ and $\frac{\partial \phi}{\partial y}=q(x, y)$ for all $(x, y) \in D$ 
\item $\frac{\partial q}{\partial x}=\frac{\partial p}{\partial y}$ holds at all points of $D$ 
\item $\oint_{C} \vec{V} \cdot d \vec{r}=0$ for every piecewise smooth closed curve $C$ in $D$ 
\item The differential $p d x+q d y$ is exact in $D$
\end{enumerate}
\hfill [JAM 2009]


 \item  Let $f, g:[-1,1] \rightarrow \mathbb{R}, f(x)=x^{3}, g(x)=x^{2}|x|$. Then 
\begin{enumerate}
   
\item $f$ and $g$ are linearly independent on $[-1,1]$ 
\item $f$ and $g$ are linearly dependent on $[-1,1]$ 
\item $f(x) g^{\prime}(x)-f^{\prime}(x) g(x)$ is NOT identically zero on $[-1,1]$ 
\item there exist continuous functions $p(x)$ and $q(x)$ such that $f$ and $g$ satisfy $y^{\prime \prime}+p y^{\prime}+q y=0$ on $[-1,1]$
 
\end{enumerate}
\hfill [JAM 2009]


\item  The value of $c$ for which there exists a twice differentiable vector field $\vec{F}$ with $\operatorname{curl} \vec{F}=2 x \hat{i}-7 y \hat{j}+c z \hat{k}$ is 
\begin{enumerate}
    
\item 0 
\item 2 
\item 5 
\item 7

\end{enumerate}
\hfill [JAM 2009]


\item  Container A contains $100 \mathrm{cc}$ of milk and container B contains $100 \mathrm{cc}$ of water. $5 \mathrm{cc}$ of the liquid in $\mathrm{A}$ is transferred to $\mathrm{B}$, the mixture is thoroughly stirred and $5 \mathrm{cc}$ of the mixture in B is transferred back into A. Each such two-way transfer is called a dilution. Let $a_{n}$ be the percentage of water in container A after $n$ such dilutions, with the understanding that $a_{0}=0$. 
\begin{enumerate}
    \item  Prove that $a_{1}=\frac{100}{21}$ and that, in general, $a_{n}=\frac{100}{21}+\frac{19}{21} a_{n-1}$ for $n=1,2,3, \ldots$ 
\item  Using (a) prove that $a_{n}=50\left[1-\left(\frac{19}{21}\right)^{n}\right]$ for $n=1,2,3, \ldots$. 
Find $\lim _{n \rightarrow \infty} a_{n}$ and explain why the answer is intuitively obvious.
\end{enumerate}
\hfill [JAM 2009]




\item

\begin{enumerate}
    \item  Let $f: \mathrm{N} \times \mathrm{N} \rightarrow \mathrm{R}$ be a non-negative function. Assume that for every $m \in \mathrm{N}$, the series $\sum_{n=1}^{\infty} f(m, n)$ is convergent and has sum $a_{m}$ and further that the series $\sum_{m=1}^{\infty} a_{m}$ is also convergent and has sum $L$. Prove that for every $n$, the series $\sum_{m=1}^{\infty} f(m, n)$ is convergent and if we denote its sum by $b_{n}$ then the series $\sum_{n=1}^{\infty} b_{n}$ is also convergent and has sum $L$.

\item  Define $f: \mathrm{N} \times \mathrm{N} \rightarrow \mathrm{R}$ by
$$
f(m, n)=\left\{\begin{array}{lll}
0 & \text { if } n>m, \\
\frac{-1}{2^{m-n}} & \text { if } n<m, \\
1 & \text { if } n=m .
\end{array}\right.
$$
Show that $\sum_{m=1}^{\infty} \sum_{n=1}^{\infty} f(m, n)=2$ and $\sum_{n=1}^{\infty} \sum_{m=1}^{\infty} f(m, n)=0$.
\end{enumerate}
\hfill [JAM 2009]



\item  
\begin{enumerate}
   \item Evaluate $\iint_{R} \cos \left(\max \left\{x^{3}, y^{3 / 2}\right\}\right) d x d y$, where $R=[0,1] \times[0,1]$.

\item  Let $S=\sqrt{1}+\sqrt{2}+\sqrt{3}+\ldots .+\sqrt{10000}$ and $I=\int_{n}^{10000} \sqrt{x} d x$. Show that $I \leq S \leq I+100$.(6)
 
\end{enumerate}
\hfill [JAM 2009]


\item  Let $D=\{(x, y): x \geq 0, y \geq 0\}$. Let $f(x, y)=\left(x^{2}+y^{2}\right) e^{-x-y}$ for $(x, y) \in D$.

Prove that $f$ attains its maximum on $D$ at two boundary points.

Deduce that $\frac{x^{2}+y^{2}}{4} \leq e^{x+y-2}$ for all $x \geq 0, y \geq 0$.
\hfill [JAM 2009]


\item  
\begin{enumerate}
    
\item Let $a_{1}, b_{1}, a_{2}, b_{2} \in \mathrm{R}$. Show that the condition $a_{2} b_{1}>0$ is sufficient but not necessary for the system
$$
\begin{aligned}
&\frac{d x}{d t}=a_{1} x+b_{1} y, \\
&\frac{d y}{d t}=a_{2} x+b_{2} y,
\end{aligned}
$$
to have two linearly independent solutions of the form $x=c_{1} e^{\lambda_{1} t}, y=d_{1} e^{\lambda_{1} t}$ and $x=c_{2} e^{\lambda_{2} t}, y=d_{2} e^{\lambda_{2} t}$ with $c_{1}, d_{1}, c_{2}, d_{2}, \lambda_{1}, \lambda_{2} \in \mathbb{R}$.

\item  Show that the differential equation representing the family of all straight lines which have an intercept of constant length $L$ between the coordinate axes is
$$
x \frac{d y}{d x}-y=\frac{L \frac{d y}{d x}}{\sqrt{1+\left(\frac{d y}{d x}\right)^{2}}} .
$$
\end{enumerate}
\hfill [JAM 2009]


\item  Let $A, B, k>0$. Solve the initial value problem

Also show that
$$
\frac{d y}{d x}-A y+B y^{3}=0, x>0, y(0)=k .
$$

\begin{enumerate}
    \item  if $k<\sqrt{\frac{A}{B}}$, then the solution $y(x)$ is monotonically increasing on $(0, \infty)$ and tends to $\sqrt{\frac{A}{B}}$ as $x \rightarrow \infty$

\item  if $k>\sqrt{\frac{A}{B}}$, then the solution $y(x)$ is monotonically decreasing on $(0, \infty)$ and tends to $\sqrt{\frac{A}{B}}$ as $x \rightarrow \infty$.
\end{enumerate}
\hfill [JAM 2009]



\item  

\begin{enumerate}
   \item Evaluate
$$
\int_{C}\left(3 y^{2}+2 z^{2}\right) d x+(6 x-10 z) y d y+\left(4 x z-5 y^{2}\right) d z
$$
along the portion from $(1,0,1)$ to $(3,4,5)$ of the curve $C$, which is the intersection of the two surfaces $z^{2}=x^{2}+y^{2}$ and $z=y+1$.

a

\item  A particle moves counterclockwise along the curve $3 x^{2}+y^{2}=3$ from $(1,0)$ to a point $P$, under the action of the force $\vec{F}(x, y)=\frac{x}{y} \hat{i}+\frac{y}{x} \hat{j}$. Prove that there are two possible locations of $P$ such that the work done by $\vec{F}$ is 1 . 
\end{enumerate}
\hfill [JAM 2009]


\item  Verify Stokes's theorem for the hemisphere $x^{2}+y^{2}+z^{2}=9, z \geq 0$ and the vector field $\vec{F}=\left(z^{2}-y\right) \hat{i}+(x-2 y z) \hat{j}+\left(2 x z-y^{2}\right) \hat{k}$.
\hfill [JAM 2009]

\item  

\begin{enumerate}
    
\item  Let $T: \mathbb{R}^{3} \rightarrow \mathbb{R}^{2}$ be the linear transformation defined by $T(x, y, z)=(x+2 y, x-z)$. Let $N(T)$ be the null space of $T$ and $W=\left\{\vec{v} \in \mathbb{R}^{3} \mid \vec{v} \cdot \vec{u}=0\right.$, for all $\left.\vec{u} \in N(T)\right\}$. Find a linear transformation $S: \mathbb{R}^{2} \rightarrow W$ such that $T S=I$, where $I$ is the identity transformation on $\mathrm{R}^{2}$.

\item  Suppose $A$ is a real square matrix of odd order such that $A+A^{T}=0$. Prove that $A$ is singular. 
\end{enumerate}
\hfill [JAM 2009]




\item  
\begin{enumerate}
    \item  Find all pairs $(a, b)$ of real numbers for which the system of equations
$$
\begin{aligned}
x+3 y &=1 \\
4 x+a y+z &=0 \\
2 x+3 z &=b
\end{aligned}
$$
has (i) a unique solution, (ii) infinitely many solutions, (iii) no solution.

\item  Let $A$ be an $n \times n$ matrix such that $A^{n}=0$ and $A^{n-1} \neq 0$. Show that there exists a vector $v \in \mathbb{R}^{\mathrm{n}}$ such that $\left\{v, A v, \ldots, A^{n-1} v\right\}$ forms a basis for $\mathbb{R}^{\mathrm{n}}$.

\end{enumerate}
\hfill [JAM 2009]


\item 
\begin{enumerate}
    
\item  In which of the following pairs are the two groups isomorphic to each other? Justify your answers.

(i) $\mathrm{R} / \mathrm{Z}$ and $S^{1}$, where $\mathrm{R}$ is the additive group of real numbers and $S^{1}=\{z \in \mathbb{C}:|z|=1\}$ under complex multiplication.

(ii) $(\mathrm{Z},+)$ and $(\mathrm{Q},+)$.



\item Prove or disprove that if $G$ is a finite abelian group of order $n$, and 
integer which divides $n$, then $G$ has at most one subgroup of order $k$.
\end{enumerate}
\hfill [JAM 2009]


\item  Let $I$ and $J$ be ideals of a ring $R$. Let $I J$ be the set of all possible sums $\sum_{i=1}^{n} a_{i} b_{i}$, where $a_{i} \in I, b_{i} \in J$ for $i=1,2, \ldots, n$ and $n \in \mathrm{N}$. 
\begin{enumerate}
    \item  Prove that $I J$ is an ideal of $R$ and $I J \subseteq I \cap J$. 
\item  Is it true that $I J=I \cap J$ ? Justify your answer.
\end{enumerate}
\hfill [JAM 2009]




\item  A sequence $\left\{f_{n}\right\}$ of functions defined on an interval $I$ is said to be uniformly bounded on $I$ if there exists some $M$ such that $\left|f_{n}(x)\right| \leq M$ for all $x \in I$ and for all $n \in \mathrm{N}$.


\begin{enumerate}
    
\item  Prove that if a sequence of functions $\left\{f_{n}\right\}$ converges to a function $f$ on $I$ and $\left\{f_{n}\right\}$ is uniformly bounded on $I$, then $f$ is bounded on $I$.

\item  Suppose the sequences $\left\{f_{n}\right\}$ and $\left\{g_{n}\right\}$ of functions converge uniformly to $f$ and $g$ respectively on $I$ and both are uniformly bounded on $I$. Prove that the product sequence $\left\{f_{n} g_{n}\right\}$ converges to $f g$ uniformly on $I$. Show by an example that this may fail if only one of $\left\{f_{n}\right\}$ and $\left\{g_{n}\right\}$ is uniformly bounded on $I$.
\end{enumerate}
\hfill [JAM 2009]


\item  
\begin{enumerate}
    
\item Prove that if $f$ is a real-valued function which is uniformly continuous on an interval $(a, b)$, then $f$ is bounded on $(a, b)$.

\item  Let $f$ be a differentiable function on an interval $(a, b)$. Assume that $f^{\prime}$ is bounded on $(a, b)$. Prove that $f$ is uniformly continuous on $(a, b)$.
\end{enumerate}
\hfill [JAM 2009]




\end{enumerate}









\newpage
\section{JAM 2010}

\begin{enumerate}


\item  Which of the following conditions does NOT ensure the convergence of a real sequence $\left\{a_{n}\right\}$ ? 

\begin{enumerate}
    
\item $\left|a_{n}-a_{n+1}\right| \rightarrow 0$ as $n \rightarrow \infty$ 
\item $\sum_{n=1}^{\infty}\left|a_{n}-a_{n+1}\right|$ is convergent 
\item $\sum_{n=1}^{\infty} n a_{n}$ is convergent 
\item The sequences $\left\{a_{2 n}\right\},\left\{a_{2 n+1}\right\}$ and $\left\{a_{3 n}\right\}$ are convergent
\end{enumerate}
\hfill [JAM 2010]




\item  The value of $\iint_{G} \frac{\ln \left(x^{2}+y^{2}\right)}{x^{2}+y^{2}} d x d y$, where $G=\left\{(x, y) \in \mathbf{R}^{2}: 1 \leq x^{2}+y^{2} \leq e^{2}\right\}$, is 

\begin{enumerate}
    
\item $\pi$ 
\item $ 2 \pi$ 
\item $3 \pi$ 
\item $4 \pi$

\end{enumerate}
\hfill [JAM 2010]



\item  The number of elements of $S_{5}$ (the symmetric group on 5 letters) which are their own inverses equals 

\begin{enumerate}
    \item 10 
\item 11 
\item $ 25$ 
\item 26

\end{enumerate}
\hfill [JAM 2010]




\item  Let $S$ be an infinite subset of $\mathbf{R}$ such that $S \cap \mathbf{Q}=\phi$. Which of the following statements is true? 

\begin{enumerate}
    \item $S$ must have a limit point which belongs 10 
\item S must have a limit point which belongs $1 \mathbf{R} \backslash \mathbf{Q}$ 
\item $S$ cannot be a closed set in $\mathbf{R}$ 
\item $\mathbf{R} \backslash S$ must have a limit point which belongs to $S$
\end{enumerate}
\hfill [JAM 2010]





\item  Let $f:(1,4) \rightarrow \mathbf{R}$ be a uniformly continuous function and let $\left\{a_{n}\right\}$ be a Cauchy sequence in $(1,2)$. Let $x_{n}=\alpha_{n}^{2} f\left(a_{n}^{2}\right)$ and $y_{n}=\frac{1}{1+a_{n}^{2}} f\left(a_{n}^{2}\right)$, for all $n \in \mathbf{N}$. Which of the following statements is true? 

\begin{enumerate}
    
\item Both $\left\{x_{n}\right\}$ and $\left\{y_{n}\right\}$ must be Cauchy sequences in $\mathbf{R}$ 
\item $\left\{x_{n}\right\}$ must be a Cauchy sequence in $\mathbf{R}$ but $\left\{y_{n}\right\}$ need not be a Cauchy sequence in $\mathbf{R}$ 
\item $\left\{y_{n}\right\}$ must be a Cauchy sequence in $R$ but $\left\{x_{n}\right\}$ need not be a Cauchy sequence in $R$ 
\item Neither $\left\{x_{n}\right\}$ nor $\left\{y_{n}\right\}$ needs to be a Cauchy sequence in $\mathbf{R}$
\end{enumerate}
\hfill [JAM 2010]




\item  Let $\vec{F}=2 x y z e^{x^{2}} \hat{i}+z e^{x^{2}} \hat{j}+y e^{x^{2}} \hat{k}$ be the gradient of a scalar function. The value of $\int_{L} \vec{F} \cdot \overrightarrow{d r}$ along the oriented path $L$ from $(0,0,0)$ to $(1,0,2)$ and then to $(1,1,2)$ is 

\begin{enumerate} 
\item 0 
\item $2 e$ 
\item $e$ 
\item $e^{2}$
\end{enumerate}
\hfill [JAM 2010]




\item  Let $\vec{F}=x y \hat{i}+y \hat{j}-y z \hat{k}$ denote the force field on a particle traversing the path $L$ from $(0,0,0)$ to $(1,1,1)$ along the curve of intersection of the cylinder $y=x^{2}$ and the plane $z=x$. The work done by $F^{\prime}$ is 

\begin{enumerate}
    \item 0 
\item $\frac{1}{4}$ 
\item $\frac{1}{2}$ 
\item 1 
\end{enumerate}
\hfill [JAM 2010]



\item  Let $\mathbf{R}[X]$ be the ring of real polynomials iu the variable $X$. The number of ideals in the quotient ring $\mathbf{R}[X] /\left(X^{2}-3 X+2\right)$ is 

\begin{enumerate}
    
\item 2 
\item 3 
\item 4 
\item 6
\end{enumerate}
\hfill [JAM 2010]




\item  Consider the differential equation $\frac{d y}{d x}=a y-b y^{2}$, where $a, b>0$ and $y(0)=y_{0}$. As $x \rightarrow+\infty$, the solution $y(x)$ tends to 

\begin{enumerate}
    \item 0 
\item $\frac{a}{b}$ 
\item $\frac{b}{a}$ 
\item $y_{0}$
\end{enumerate}
\hfill [JAM 2010]





\item  Consider the differential equation $(x+y+1) d x+(2 x+2 y+1) d y=0$. Which of the following statements is true? 

\begin{enumerate}
    
\item The differential equation is linear 
\item The differential equation is exact 
\item $e^{x+y}$ is an integrating factor of the differential equation 
\item A suitable substitution transforms the differentiable equation to the variables separable form
\end{enumerate}
\hfill [JAM 2010]




\item  Let $T: \mathbf{R}^{2} \rightarrow \mathbf{R}^{2}$ be a linear transformation such that $T((1,2))=(2,3)$ and $T((0,1))=(1,4)$. Then $T((5,6))$ is 

\begin{enumerate}
\item $(6,-1)$ 
\item $(-6,1)$ 
\item $(-1,6)$ 
\item $(1,-6)$
\end{enumerate}
\hfill [JAM 2010]




\item  The number of $2 \times 2$ matrices over $\mathbf{Z}_{3}$ (the field with three elements) with determinant 1 is 

\begin{enumerate}
    \item 24 
\item 60 
\item 20 
\item 30
\end{enumerate}
\hfill [JAM 2010]





\item  The radius of convergence of the power series $\sum_{n=0}^{n} a_{n} 2^{n^{2}}$, where $a_{0}=1, a_{n}=3^{-n} a_{n-1}$ for $n \in \mathbf{N}$, is 

\begin{enumerate}
    
\item 0 
\item $\sqrt{3}$ 
\item 3 
\item $\quad \infty$
\end{enumerate}
\hfill [JAM 2010]




\item  Let $T: \mathbf{R}^{3} \rightarrow \mathbf{R}^{3}$ be the linear transformation whose matrix with respect to the standard basis $\left\{e_{1}, e_{2}, e_{9}\right\}$ of $\mathbf{R}^{3}$ is $\left(\begin{array}{lll}0 & 0 & 1 \\ 0 & 1 & 0 \\ 1 & 0 & 0\end{array}\right)$. Then $T$ 

\begin{enumerate}
    
\item maps the subspace spanned by $e_{1}$ and $e_{2}$ into itself 
\item has distinct eigenvalues 
\item has eigenvectors that span $\mathbf{R}^{3}$ 
\item has a non-zero null space 
\end{enumerate}
\hfill [JAM 2010]


\item  Let $T: \mathbf{R}^{3} \rightarrow \mathbf{R}^{9}$ be the linear transformation whose matrix with respect to the standard basis of $\mathbf{R}^{3}$ is $\left(\begin{array}{ccc}0 & a & b \\ -a & 0 & c \\ -b & -c & 0\end{array}\right)$, where $a, b, c$ are real numbers not all zero. Then $T$ 

\begin{enumerate}
    
\item is one-to-one 
\item is onto 
\item does not map any line through the origin onto itself 
\item has rank 1
\end{enumerate}
\hfill [JAM 2010]




\item 

\begin{enumerate}
     \item  Obtain the general solution of the following system of differential equations:
$$
\begin{aligned}
&\frac{d x}{d t}=x+2 y \\
&\frac{d y}{d t}=4 x-y+e^{y t}
\end{aligned}
$$
\item  Find the curve passing through $\left(\frac{1}{2}, 0\right)$ and having slope at $(x, y)$ given by the differential equatiou $2\left(1+y^{2}\right) d x+\left(2 x-\tan ^{-1} y\right) d y=0$
\end{enumerate}
\hfill [JAM 2010]




\item  

\begin{enumerate}
    \item  Find the volume of the region in the first octant bounded by the surfaces $x=0, y=x, y=2-x^{2}, z=0$ and $z=x^{2}$.

\item Suppose $f: \mathbf{R} \rightarrow \mathbf{R}$ is a non-constant continuous function satisfying $f(x+y)=f(x) f(y)$ for all $x, y \in \mathbf{R}$.


\begin{enumerate}
    \item  Show that $f(x) \neq 0$ for all $x \in \mathbf{R}$.

\item  Show that $f(x)>0$ for all $x \in \mathbf{R}$.

\item  Show that there exists $\beta \in \mathbf{R}$ such that $f(x)=\beta^{x}$ for all $x \in \mathbf{R}$.
\end{enumerate}

\end{enumerate}
\hfill [JAM 2010]



\item  

\begin{enumerate}
    \item Let $f(x)$ and $g(x)$ be real valued functions continuous in $\lfloor\alpha, b]$, differentiable in $(a, b)$ and let $\xi^{\prime}(x) \neq 0$ for all $x \in(a, b)$. Show that there exists $c \in(a, b)$ such that $\frac{f(c)-f(a)}{g(b)-g(c)}=\frac{f^{\prime}(c)}{g^{\prime}(c)}$

\item  Let $0<\lambda<4$ and let $\left\{a_{n}\right\}$ be a sequence of positive real numbers satisfying $a_{n+1}=\lambda a_{n}^{2}\left(1-a_{n}\right)$ for $n \in \mathbf{N}$. Prove that $\lim _{n \rightarrow-\infty} a_{n}$ exists and determine this limit.
\end{enumerate}
\hfill [JAM 2010]




\item  Let $G$ be an open subset of $\mathbf{R}$.


\begin{enumerate}
    \item  If $0 \notin G$, then show that $H=\{x y: x, y \in G\}$ is an open subset of $\mathbf{R}$.

\item  If $0 \in G$ and if $x+y \in G$ for all $x, y \in G$, then show that $G=\mathbf{R}$.
\end{enumerate}
\hfill [JAM 2010]





\item  Let $p(x)$ be a non-constant polynomial with real coefficients such that $p(x) \neq 0$ for all $x \in \mathbf{R}$. Define $f(x)=\frac{1}{p(x)}$ for all $x \in \mathbf{R}$. Prove that


\begin{enumerate}
    \item  for each $\varepsilon>0$, there exists $\alpha>0$ such that $|f(x)|<\varepsilon$ for all $x \in \mathbf{R}$ satisfying $|x|>\alpha$, and

\item  $f: \mathbf{R} \rightarrow \mathbf{R}$ is a uniformly continuous function. 
\end{enumerate}
\hfill [JAM 2010]






\item 

\begin{enumerate}
    Let $M(k)$ and $m(k)$ denote respectively the absolute maximum and the absolute minimum values of $x^{3}+9 x^{2}-21 x+k$ in the closed interval $[-10,2]$. Find all the real values of $k$ for which $|M(k)|=|m(k)|$.

\item Let $\alpha_{1}=0, \beta_{1}=1 ; \alpha_{2}=1, \beta_{2}=1$, and for $n \geq 3$,

$\alpha_{n}=\alpha_{n-1}+2 \alpha_{n-2}$

$\beta_{n}=\beta_{n-1}+2 \beta_{n-2}$

Prove that, for $n \in \mathbf{N}$


\begin{enumerate}
\item  $\beta_{n}=2 a_{n}+(-1)^{n-1}$
\item  $\alpha_{n}+\beta_{n}=2^{n-1}$
\end{enumerate}

Deduce that $\lim _{n \rightarrow \infty} \frac{a_{\mathrm{n}} a+\beta_{n} b}{2^{n-1}}=\frac{a+2 b}{3}$ for any $a, b \in \mathbf{R}$.

\end{enumerate}
\hfill [JAM 2010]





\item  

\begin{enumerate}
    
\item  Let $f(x, y)=\alpha x^{2}+x y+\beta y^{2}, \alpha \neq 0, \beta \neq 0,4 \alpha \beta \neq 1$. Find suflicient conditions on $\alpha, \rho$ such that $(0,0)$ is


\begin{enumerate}
    \item a point of local maxima of $f(x, y)$

    \item  a point of local minima of $f(x, y)$

    \item  a saddle point of $f(x, y)$.
\end{enumerate}


\item  Find the derivative of $f(x, y, z)=7 x^{3}-x^{2} z-z^{2}+28 y$ at the point $A=(1,-1,0)$ along the unit vector $\frac{1}{7}(6 \hat{i}-2 \hat{j}+3 \hat{k})$. What is the unit vector along which $f$ decreases most rapidly at $A$ ? Also, find the rate of this decrease.

\end{enumerate}
\hfill [JAM 2010]


\item  Using $x=e^{u}$, transform the differential equation $x^{2} \frac{d^{2} y}{d x^{2}}+4 x \frac{d y}{d x}+2 y=\cos x$ to a second order differential equation with constant coeflicients. Obtein the general solntion of the transformed diflerential equation.
\hfill [JAM 2010]




\item  Let $G$ be a group and let $A(G)$ denote the set of all automorphisms of $G$, i.e., all one-to-one, onto, group homomorphisms from $G$ to $G$. An automorphism $f: G \rightarrow G$ of the form $f(x)=a x a^{-1}, x \in G$ (for some $a \in G$ ) is called an inner automorphism. Let $I(G$ ) denote the set of all inner automorphisms of $G$.


\begin{enumerate}
   
\item  Show that $A(G)$ is a group under composition of functions and that $I(G)$ is a normal subgroup of $A(G)$.

\item  Show that $I(G)$ is isomorphic to $G / Z(G)$, where $Z(G)=\{g \in G: x g=g x$ for all $x \in G\}$ is the center of $G$. 
\end{enumerate}
\hfill [JAM 2010]




\item  

\begin{enumerate}
    
\item  Give an example of a linear transformation $T: \mathbf{R}^{2} \rightarrow \mathbf{R}^{2}$ such that $T^{2}(v)=-v$ for all $v \in \mathbf{R}^{2}$.

\item  Let $V$ be a real n-dimensional vector space and let $T: V \rightarrow V$ be a linear transformation satisfying $T^{2}(v)=-v$ for all $v \in V$.



\begin{enumerate}
    \item  Show that $n$ is even.

\item  Use $T$ to make $V$ into a complex vector space such that the multiplication by complex numbers extends the multiplication by real numbers.

\item  Show that, with respect to the complex vector space structure on $V$ obtained in (ii), $T: V \rightarrow V$ is a complex linear transformation. 
\end{enumerate}
\end{enumerate}
\hfill [JAM 2010]



\item  Let $W$ be the region bounded by the planes $x=0, y=0, y=3, z=0$ and $x+2 z=6$.
Let $S$ be the boundary of this region. Using Gauss' divergence theorem, evaluate $\iint_{S} \vec{F} \cdot \hat{n} d S$, where $\vec{F}=2 x y \hat{i}+y z^{2} \hat{j}+x z \hat{k}$ and $\hat{n}$ is the outward unit normal vector to $S$.
\hfill [JAM 2010]

\item  

\begin{enumerate}
    \item  Using Stokes' theorem evaluate the line integral $\int_{L}(y \hat{i}+z \hat{j}+x \hat{k}) \cdot \overrightarrow{d r}$, where $L$ is the intersection of $x^{2}+y^{2}+z^{2}=1$ and $x+y=0$ traversed in the clockwise direction when viewed from the point $(1,1,0)$.

\item  Change the order of integration in the integral $\int_{0}^{1} \int_{x-1}^{\sqrt{1-x^{2}}} f(x, y) d y d x$.
\end{enumerate}
\hfill [JAM 2010]




\item  In a group $G, x \in G$ is said ta be conjugate to $y \in G$, written $x \sim y$, if there exists $z \in G$ such that $x=z y z^{-1}$.


\begin{enumerate}
    \item  Show that $\sim$ is an equivalence relation on $G$. Show that a subgroup $N$ of $G$ is a normal subgroup of $G$ if and only if $N$ is a union of equivalence classes of $~$.

\item  Consider the group of all non-singular $3 \times 3$ real matrices under matrix multiplication. Show that $\left[\begin{array}{lll}1 & 0 & 0 \\ 1 & 3 & 0 \\ 1 & 2 & 1\end{array}\right] \sim\left[\begin{array}{lll}3 & 0 & 4 \\ 0 & 1 & 0 \\ 0 & 0 & 1\end{array}\right]$ (i.e., the two matrices are conjugate).

\end{enumerate}
\hfill [JAM 2010]




\item  Let $S$ denote the commutative ring of all continuous real-valued functions on $[0,1]$, under pointwise addition and multiplication. For $a \in[0,1]$, let $\left.M_{a}=|f \in S| f(a)=0\right\}$. 


\begin{enumerate}
    
\item  Show that $M_{a}$ is an ideal in $S$. 
\item  Show that $M_{\text {at }}$ is a maximal ideal in $S$.


\end{enumerate}
\hfill [JAM 2010]



\end{enumerate}

\newpage

\section{JAM 2011}

\begin{enumerate}
    
\item  Let $a_{n}=\sum_{k=1}^{n} \frac{n}{n^{2}+k}$, for $n \in \mathbb{N}$. Then the sequence $\left\{a_{n}\right\}$ is 
\begin{enumerate}
    \item Convergent 
\item Diverges to $\infty$ 
\item Bounded but not convergent 
\item Neither bounded nor diverges to $\infty$
\end{enumerate}
\hfill [JAM 2011]




\item  The number of real roots of the equation $x^{3}+x-1=0$ is 

\begin{enumerate}
\item 0 
\item 1 
\item 2 
\item 3
\end{enumerate}
\hfill [JAM 2011]



\item  The value of $\lim _{n \rightarrow \infty} \sum_{k=1}^{n} \frac{1}{\sqrt{n^{2}+k n}}$ is 
\begin{enumerate}
    \item $2(\sqrt{2}-1)$ 
\item $2 \sqrt{2}-1$ 
\item $2-\sqrt{2}$ 
\item $\frac{1}{2}(\sqrt{2}-1)$
\end{enumerate}
\hfill [JAM 2011]




\item  Let $V$ be the region bounded by the planes $x=0, x=2, y=0, z=0$ and $y+z=1$. Then the value of the integral $\iiint_{V} y d x d y d z$ is 
\begin{enumerate}
    \item $\frac{1}{2}$ 
\item $\frac{4}{3}$ 
\item $1 $
\item $\frac{1}{3}$
\end{enumerate}
\hfill [JAM 2011]




\item  The solution $y(x)$ of the differential equation $\frac{d^{2} y}{d x^{2}}+4 \frac{d y}{d x}+4 y=0$ satisfying the conditions $y(0)=4, \frac{d y}{d x}(0)=8$ is 
\begin{enumerate}
    \item $4 e^{2 x}$ 
\item $(16 x+4) e^{-2 x}$ 
\item $4 e^{-2 x}+16 x$ 
\item $4 e^{-2 x}+16 x e^{2 x}$
\end{enumerate}
\hfill [JAM 2011]




\item  If $y^{a}$ is an integrating factor of the differential equation $2 x y d x-\left(3 x^{2}-y^{2}\right) d y=0$, then the value of $a$ is 
\begin{enumerate}
    \item $-4$ 
\item $4$ 
\item $-1$ 
\item $1$

\end{enumerate}
\hfill [JAM 2011]



\item  Let $\vec{F}=a y \hat{i}+z \hat{j}+x \hat{k}$ and $C$ be the positively oriented closed curve given by $x^{2}+y^{2}=1, z=0$. If $\oint_{C} \vec{F} \cdot d \vec{r}=\pi$, then the value of $a$ is 
\begin{enumerate}
    \item $-1$ 
\item $0$ 
\item $\frac{1}{2}$ 
\item $1$
\end{enumerate}
\hfill [JAM 2011]




\item  Consider the vector field $\vec{F}=(a x+y+a) \hat{i}+\hat{j}-(x+y) \hat{k}$, where $a$ is a constant. If $\vec{F} \cdot \operatorname{cur} 1 \vec{F}=0$, then the value of $a$ is 
\begin{enumerate}
    \item $-1$ 
\item 0 
\item 1 
\item $\frac{3}{2}$
\end{enumerate}
\hfill [JAM 2011]


 \item  Let $G$ denote the group of all $2 \times 2$ invertible matrices with entries from $\mathbb{R}$. Let
$$
H_{1}=\{A \in G: \operatorname{det}(A)=1\} \text { and } H_{2}=\{A \in G: A \text { is upper triangular }\} \text {. }
$$
Consider the following statements:

$P: H_{1}$ is a normal subgroup of $G$

$Q: H_{2}$ is a normal subgroup of $G$.

Then
\begin{enumerate}
    \item Both $P$ and $Q$ are true 
\item $P$ is true and $Q$ is false 
\item $P$ is false and $Q$ is true 
\item Both $P$ and $Q$ are false
\end{enumerate}
\hfill [JAM 2011]

 


\item  For $n \in \mathbb{N}$, let $n \mathbb{Z}=\{n k: k \in \mathbb{Z}\}$. Then the number of units of $\mathbb{Z} / 11 \mathbb{Z}$ and $\mathbb{Z} / 12 \mathbb{Z}$, respectively, are 
\begin{enumerate}
    \item 11,12 
\item 10,11 
\item 10,4 
\item 10,8
\end{enumerate}
\hfill [JAM 2011]




\item  Let $A$ be a $3 \times 3$ matrix with trace $(A)=3$ and $\operatorname{det}(A)=2$. If 1 is an eigenvalue of $A$, then the eigenvalues of the matrix $A^{2}-2 I$ are 
\begin{enumerate}
    \item $1,2(i-1),-2(i+1)$ 
\item $-1,2(i-1), 2(i+1)$ 
\item $1,2(i+1),-2(i+1)$ 
\item $-1,2(i-1),-2(i+1)$

\end{enumerate}
\hfill [JAM 2011]



\item  Let $T: \mathbb{R}^{n} \rightarrow \mathbb{R}^{n}$ be a linear transformation, where $n \geq 2$. For $k \leq n$, let

Then
$$
E=\left\{v_{1}, v_{2}, \ldots, v_{k}\right\} \subseteq \mathbb{R}^{n} \text { and } F=\left\{T v_{1}, T v_{2}, \ldots, T v_{k}\right\}
$$
\begin{enumerate}
    \item If $E$ is linearly independent, then $F$ is linearly independent 
\item If $F$ is linearly independent, then $E$ is linearly independent 
\item If $E$ is linearly independent, then $F$ is linearly dependent 
\item If $F$ is linearly independent, then $E$ is linearly dependent

\end{enumerate}
\hfill [JAM 2011]



\item  For $n \neq m$, let $T_{1}: \mathbb{R}^{n} \rightarrow \mathbb{R}^{m}$ and $T_{2}: \mathbb{R}^{m} \rightarrow \mathbb{R}^{n}$ be linear transformations such that $T_{1} T_{2}$ is bijective. Then
\begin{enumerate}
     
\item $\operatorname{rank}\left(T_{1}\right)=n$ and $\operatorname{rank}\left(T_{2}\right)=m$ 
\item rank $\left(T_{1}\right)=m$ and $\operatorname{rank}\left(T_{2}\right)=n$ 
\item $\operatorname{rank}\left(T_{1}\right)=n$ and $\operatorname{rank}\left(T_{2}\right)=n$ 
\item $\operatorname{rank}\left(T_{1}\right)=m$ and $\operatorname{rank}\left(T_{2}\right)=m$
\end{enumerate}
\hfill [JAM 2011]



\item  The set of all $x$ at which the power series $\sum_{n=1}^{\infty} \frac{n}{(2 n+1)^{2}}(x-2)^{3 n}$ converges is 
\begin{enumerate}
    \item $[-1,1)$ 
\item $[-1,1]$ 
\item $[1,3)$ 
\item $[1,3]$ 

\end{enumerate}
\hfill [JAM 2011]

% MAINUL

\item  Consider the following subsets of $\mathbb{R}$ :
Then
$$
E=\left\{\frac{n}{n+1}: n \in \mathbb{N}\right\}, \quad F=\left\{\frac{1}{1-x}: 0 \leq x<1\right\}
$$
\begin{enumerate}
    \item Both $E$ and $F$ are closed 
\item $E$ is NOT closed and $F$ is closed 
\item $E$ is closed and $F$ is NOT closed 
\item Neither $E$ nor $F$ is closed
\end{enumerate}
\hfill [JAM 2011]




\item  
\begin{enumerate}
    \item Let $\left\{a_{n}\right\}$ be a sequence of non-negative real numbers such that $\sum_{n=1}^{\infty} a_{n}$ converges, and let $\left\{k_{n}\right\}$ be a strictly increasing sequence of positive integers. Show that
$$
\sum_{n=1}^{\infty} a_{k_{n}} \text { also converges. }
$$
\item  Suppose $f:[0,1] \rightarrow \mathbb{R}$ is differentiable and $f^{\prime}(x) \leq 1$ at every $x \in(0,1)$. If $f(0)=0$ and $f(1)=1$, show that $f(x)=x$ for all $x \in[0,1]$.
\end{enumerate}
\hfill [JAM 2011]



\item 

\begin{enumerate}
    \item Suppose $f$ is a real valued function defined on an open interval $I$ and differentiable at every $x \in I$. If $[a, b] \subset I$ and $f^{\prime}(a)<0<f^{\prime}(b)$, then show that there exists $c \in(a, b)$ such that $f(c)=\min _{a \leq x \leq b} f(x)$.

\item  Let $f:(a, b) \rightarrow \mathbb{R}$ be a twice differentiable function such that $f^{\prime \prime}$ is continuous at every point in $(a, b)$. Prove that
$$
\lim _{h \rightarrow 0} \frac{f(x+h)-2 f(x)+f(x-h)}{h^{2}}=f^{\prime \prime}(x)
$$
for every $x \in(a, b)$.
\end{enumerate}
\hfill [JAM 2011]



\item  Find all critical points of the following function and check whether the function attains maximum or minimum at each of these points:
$$
u(x, y)=x^{4}+y^{4}-2 x^{2}-2 y^{2}+4 x y, \quad(x, y) \in \mathbb{R}^{2}
$$
\hfill [JAM 2011]

\item  
\begin{enumerate}
    \item  Let $\varphi:[a, b] \rightarrow \mathbb{R}$ be differentiable and $[c, d]=\{\varphi(x): a \leq x \leq b\}$, and let $f:[c, d] \rightarrow \mathbb{R}$ be continuous. Let $g:[a, b] \rightarrow \mathbb{R}$ be defined by $g(x)=\int_{c}^{\varphi(x)} f(t) d t$ for $x \in[a, b]$. Then show that $g$ is differentiable and $g^{\prime}(x)=f(\varphi(x)) \varphi^{\prime}(x)$ for all $x \in[a, b]$

\item If $f:[0,1] \rightarrow \mathbb{R}$ is such that $\int_{0}^{\sin x} f(t) d t=\frac{\sqrt{3}}{2} x$ for all $x \in \mathbb{R}$, then find $f\left(\frac{1}{2}\right)$.
\end{enumerate}
\hfill [JAM 2011]



\item  Find the area of the surface of the solid bounded by the cone $z=3-\sqrt{x^{2}+y^{2}}$ and the paraboloid $z=1+x^{2}+y^{2}$ \item  Obtain the general solution of each of the following differential equations: 

\begin{enumerate}
    \item $y-x \frac{d y}{d x}=\frac{d y}{d x} y^{2} e^{y}$. 
\item $\frac{d y}{d x}=\frac{x+2 y+8}{2 x+y+7}$.
\end{enumerate}
\hfill [JAM 2011]




\item  

\begin{enumerate}
    \item Determine the values of $b>1$ such that the differential equation
$$
x^{2} \frac{d^{2} y}{d x^{2}}+x \frac{d y}{d x}+y=0, \quad 1<x<b
$$
satisfying the conditions $y(1)=0=y(b)$ has a nontrivial solution.

\item  Find $v(x)$ such that $y(x)=e^{4 x} v(x)$ is a particular solution of the differential equation
$$
\frac{d^{2} y}{d x^{2}}-8 \frac{d y}{d x}+16 y=\left(2 x+11 x^{10}+21 x^{20}\right) e^{4 x} .
$$
\end{enumerate}
\hfill [JAM 2011]



\item 
\begin{enumerate}
    
  \item Change the order of integration in the double integral
$$
\int_{-1}^{2}\left(\int_{-x}^{2-x^{2}} f(x, y) d y\right) d x
$$
\item  Let $\vec{F}=\left(x^{2}-x y^{2}\right) \hat{i}+y^{2} \hat{j}$. Using Green's theorem, evaluate the line integral $\int_{C} \vec{F} \cdot d \vec{r}$, where $C$ is the positively oriented closed curve which is the boundary of the region enclosed by the $x$ - axis and the semi-circle $y=\sqrt{1-x^{2}}$ in the upper half plane.

\end{enumerate}
\hfill [JAM 2011]




\item 
\begin{enumerate}
    
\item  If $\vec{F}=\left(x^{2}+y-4\right) \hat{i}+3 x y \hat{j}+\left(2 x z+z^{2}\right) \hat{k}$, then evaluate the surface integral $\iint_{S}(\nabla \times \vec{F}) \cdot \hat{n} d S$, where $S$ is the surface of the cone $z=1-\sqrt{x^{2}+y^{2}}$ lying above the $x y$-plane and $\hat{n}$ is the unit normal to $S$ making an acute angle with $\hat{k}$.

\item  Show that the series $\sum_{n=1}^{\infty} \frac{x}{\sqrt{n}\left(1+n^{p} x^{2}\right)}$ converges uniformly on $\mathbb{R}$ for $p>1$.
\end{enumerate}
\hfill [JAM 2011]


 \item  
\begin{enumerate}
    \item  Find a value of $c$ such that the following system of linear equations has no solution:
$$
\begin{gathered}
x+2 y+3 z=1,  
3 x+7 y+c z=2,  
2 x+c y+12 z=3 .
\end{gathered}
$$
\item  Let $V$ be the vector space of all polynomials with real coefficients of degree at most $n$, where $n \geq 2$. Considering elements of $V$ as functions from $\mathbb{R}$ to $\mathbb{R}$, define
$$
W=\left\{p \in V: \int_{0}^{1} p(x) d x=0\right\} .
$$
Show that $W$ is a subspace of $V$ and $\operatorname{dim}(W)=n$.

\end{enumerate}
\hfill [JAM 2011]


\item
\begin{enumerate}
     \item  Let $A$ be a $3 \times 3$ real matrix with $\operatorname{det}(A)=6$. Then find $\operatorname{det}(\operatorname{adj} A)$.

\item  Let $v_{1}$ and $v_{2}$ be non-zero vectors in $\mathbb{R}^{n}, n \geq 3$, such that $v_{2}$ is not a scalar multiple of $v_{1}$. Prove that there exists a linear transformation $T: \mathbb{R}^{n} \rightarrow \mathbb{R}^{n}$ such that $T^{3}=T, T v_{1}=v_{2}$, and $T$ has at least three distinct eigenvalues.
\end{enumerate}
\hfill [JAM 2011]

 

\item  
\begin{enumerate}
    \item  If $E$ is a subset of $\mathbb{R}$ that does not contain any of its limit points, then prove that $E$ is a countable set.


\item Let $f:(a, b) \rightarrow \mathbb{R}$ be a continuous function. If $f$ is uniformly continuous, then prove that there exists a continuous function $g:[a, b] \rightarrow \mathbb{R}$ such that $g(x)=f(x)$ for all $x \in(a, b)$.
\end{enumerate}
\hfill [JAM 2011]



\item  
\begin{enumerate}
    
\item  On $\mathbb{R}^{3}$, define a binary operation * as follows: For $(x, y, t),\left(x^{\prime}, y^{\prime}, t^{\prime}\right)$ in $\mathbb{R}^{3}$,
$$
(x, y, t)^{*}\left(x^{\prime}, y^{\prime}, t^{\prime}\right)=\left(x+x^{\prime}, y+y^{\prime}, t+t^{\prime}+\frac{1}{2}\left(x^{\prime} y-x y^{\prime}\right)\right) \text {. }
$$
Then show that $\left(\mathbb{R}^{3},{ }^{*}\right)$ is a group, and find its center.

\item  For $k \in \mathbb{N}$, let $k \mathbb{Z}=\{k n: n \in \mathbb{Z}\}$. For any $m, n \in \mathbb{N}$, show that $I=m \mathbb{Z} \cap n \mathbb{Z}$ is an ideal of $\mathbb{Z}$. Further, find the generators of $I$.
\end{enumerate}
\hfill [JAM 2011]


\item  Let $G$ be a group of order $p^{2}$, where $p$ is a prime number. Let $x \in G$. Prove that $\{y \in G: x y=y x\}=G$.
\hfill [JAM 2011]




\end{enumerate}

\newpage
\section{JAM 2012}

\begin{enumerate}
  

\item  Let $\left\{x_{n}\right\}$ be the sequence $+\sqrt{1},-\sqrt{1},+\sqrt{2},-\sqrt{2},+\sqrt{3},-\sqrt{3},+\sqrt{4},-\sqrt{4}, \ldots$.

If
$$
y_{n}=\frac{x_{1}+x_{2}+\cdots+x_{n}}{n} \text { for all } n \in \mathbb{N},
$$
then the sequence $\left\{y_{n}\right\}$ is


\begin{enumerate}
    \item monotonic 
\item NOT bounded 
\item bounded but NOT convergent 
\item convergent
\end{enumerate}
\hfill [JAM 2012]





\item  The number of distinct real roots of the equation $x^{9}+x^{7}+x^{5}+x^{3}+x+1=0$ is 


\begin{enumerate}
    
\item 1 
\item 3 
\item 5 
\item 9
\end{enumerate}
\hfill [JAM 2012]




\item  If $f: \mathbb{R}^{2} \rightarrow \mathbb{R}$ is defined by
$$
f(x, y)=\left\{\begin{array}{cc}
\frac{x^{3}}{x^{2}+y^{4}} & \text { if }(x, y) \neq(0,0),  \\
0 & \text { if }(x, y)=(0,0),
\end{array}\right.
$$
then


\begin{enumerate}
   \item $f_{x}(0,0)=0$ and $f_{y}(0,0)=0$ 
\item $f_{x}(0,0)=1$ and $f_{y}(0,0)=0$ 
\item $f_{x}(0,0)=0$ and $f_{y}(0,0)=1$ 
\item $f_{x}(0,0)=1$ and $f_{y}(0,0)=1$ 
\end{enumerate}
\hfill [JAM 2012]





\item  The value of $\int_{z=0}^{1} \int_{y=0}^{z} \int_{x=0}^{y} x y^{2} z^{3} d x d y d z$ is 


\begin{enumerate}
    
\item $\frac{1}{90}$ 
\item $\frac{1}{50}$ 
\item $\frac{1}{45}$ 
\item $\frac{1}{10}$
\end{enumerate}
\hfill [JAM 2012]




\item  The differential equation $\left(1+x^{2} y^{3}+\alpha x^{2} y^{2}\right) d x+\left(2+x^{3} y^{2}+x^{3} y\right) d y=0$ is exact if $\alpha$ equals 


\begin{enumerate}
    
\item $\frac{1}{2}$ 
\item $\frac{3}{2}$ 
\item 2 
\item 3
\end{enumerate}
\hfill [JAM 2012]




\item  An integrating factor for the differential equation $\left(2 x y+3 x^{2} y+6 y^{3}\right) d x+\left(x^{2}+6 y^{2}\right) d y=0$ is 


\begin{enumerate}
    \item $x^{3}$ 
\item $y^{3}$ 
\item $e^{3 x}$ 
\item $e^{3 y}$
\end{enumerate}
\hfill [JAM 2012]




\item  For $c>0$, if $a \hat{i}+b \hat{j}+c \hat{k}$ is the unit normal vector at $(1,1, \sqrt{2})$ to the cone $z=\sqrt{x^{2}+y^{2}}$, then 


\begin{enumerate}
    \item $\quad a^{2}+b^{2}-c^{2}=0$ 
\item $-a^{2}+b^{2}+c^{2}=0$ 
\item $\quad a^{2}-b^{2}+c^{2}=0$ 
\item $a^{2}+b^{2}+c^{2}=0$
\end{enumerate}
\hfill [JAM 2012]





\item  Consider the quotient group $\mathbb{Q} / \mathbb{Z}$ of the additive group of rational numbers. The order of the element $\frac{2}{3}+\mathbb{Z}$ in $\mathbb{Q} / \mathbb{Z}$ is 


\begin{enumerate}
    
\item 2 
\item 3 
\item 5 
\item 6
\end{enumerate}
\hfill [JAM 2012]


 \item  Which one of the following is TRUE ? 
 

\begin{enumerate}
    \item The characteristic of the ring $6 \mathbb{Z}$ is 6 
\item The ring $6 \mathbb{Z}$ has a zero divisor 
\item The characteristic of the ring $(\mathbb{Z} / 6 \mathbb{Z}) \times 6 \mathbb{Z}$ is zero 
\item The ring $6 \mathbb{Z} \times 6 \mathbb{Z}$ is an integral domain
\end{enumerate}
\hfill [JAM 2012]





\item  Let $W$ be a vector space over $\mathbb{R}$ and let $T: \mathbb{R}^{6} \rightarrow W$ be a linear transformation such that $S=\left\{T e_{2}, T e_{4}, T e_{6}\right\}$ spans $W$. Which one of the following must be TRUE ? 

\begin{enumerate}
\item $S$ is a basis of $W$ 
\item $T\left(\mathbb{R}^{6}\right) \neq W$ 
\item $\left\{T e_{1}, T e_{3}, T e_{5}\right\}$ spans $W$ 
\item $\operatorname{ker}(T)$ contains more than one element

\end{enumerate}
\hfill [JAM 2012]





\item  Consider the following subspace of $\mathbb{R}^{3}$ :
$$
W=\left\{(x, y, z) \in \mathbb{R}^{3} \mid 2 x+2 y+z=0,3 x+3 y-2 z=0, x+y-3 z=0\right\} .
$$
The dimension of $W$ is 


\begin{enumerate}
    \item 0 
\item 1 
\item 2 
\item 3
\end{enumerate}
\hfill [JAM 2012]





\item  Let $P$ be a $4 \times 4$ matrix whose determinant is 10 . The determinant of the matrix $-3 P$ 
\begin{enumerate}
  \item $-810$ 
\item $-30$ 
\item 30 
\item 810  
\end{enumerate}
\hfill [JAM 2012]





\item  If the power series $\sum_{n=0}^{\infty} a_{n} x^{n}$ converges for $x=3$, then the series $\sum_{n=0}^{\infty} a_{n} x^{n}$ 

\begin{enumerate}
    \item converges absolutely for $x=-2$ 
\item converges but not absolutely for $x=-1$ 
\item converges but not absolutely for $x=1$ 
\item diverges for $x=-2$
\end{enumerate}
\hfill [JAM 2012]





\item  If $Y=\left\{\frac{x}{1+|x|} \mid x \in \mathbb{R}\right\}$, then the set of all limit points of $Y$ is 

\begin{enumerate}
    \item $(-1,1)$ 
\item $\quad(-1,1]$ 
\item $[0,1]$ 
\item $[-1,1]$
\end{enumerate}
\hfill [JAM 2012]





\item  If $C$ is a smooth curve in $\mathbb{R}^{3}$ from $(0,0,0)$ to $(2,1,-1)$, then the value of
$$
\int_{C}(2 x y+z) d x+\left(z+x^{2}\right) d y+(x+y) d z
$$
is 

\begin{enumerate}
    \item $-1$ 
\item 0 
\item 1 
\item 2
\end{enumerate}
\hfill [JAM 2012]







\item 

\begin{enumerate}
     \item Examine whether the following series is convergent:
$$
\sum_{n=1}^{\infty} \frac{n !}{1 \cdot 3 \cdot 5 \cdots(2 n-1)} \text {. }
$$
\item  For each $x \in \mathbb{R}$, let $[x]$ denote the greatest integer less than or equal to $x$. Further, for a fixed $\beta \in(0,1)$, define $a_{n}=\frac{1}{n}[n \beta]+n^{2} \beta^{n}$ for all $n \in \mathbb{N}$. Show that the sequence $\left\{a_{n}\right\}$ converges to $\beta$.
\end{enumerate}
\hfill [JAM 2012]




\item  


\begin{enumerate}
    \item Evaluate $\lim _{x \rightarrow 0} \frac{\int_{0}^{x^{2}} \sqrt{4+t^{3}} d t}{x^{2}}$.

\item For $a, b \in \mathbb{R}$ with $a<b$, let $f:[a, b] \rightarrow \mathbb{R}$ be continuous on $[a, b]$ and twice differentiable on $(a, b)$. Further, assume that the graph of $f$ intersects the straight line segment joining the points $(a, f(a))$ and $(b, f(b))$ at a point $(c, f(c))$ for $a<c<b$. Show that there exists a real number $\xi \in(a, b)$ such that $f^{\prime \prime}(\xi)=0$.

\end{enumerate}
\hfill [JAM 2012]




\item 


\begin{enumerate}
    \item Show that the point $(0,0)$ is neither a point of local minimum nor a point of local maximum for the function $f: \mathbb{R}^{2} \rightarrow \mathbb{R}$ given by $f(x, y)=3 x^{4}-4 x^{2} y+y^{2}$ for $(x, y) \in \mathbb{R}^{2}$

\item Find all the critical points of the function $f: \mathbb{R}^{2} \rightarrow \mathbb{R}$ given by $f(x, y)=x^{3}+y^{3}-3 x-12 y+40$ for $(x, y) \in \mathbb{R}^{2}$. Also, examine whether the function $f$ attains a local maximum or a local minimum at each of these critical points.
\end{enumerate}
\hfill [JAM 2012]





\item  


\begin{enumerate}
    \item Evaluate $\int_{x=0}^{4} \int_{y=\sqrt{4-x}}^{2} e^{y^{3}} d y d x$.

\item Using multiple integral, find the volume of the solid region in $\mathbb{R}^{3}$ bounded above by the hemisphere $z=1+\sqrt{1-x^{2}-y^{2}}$ and bounded below by the cone $z=\sqrt{x^{2}+y^{2}}$.

\end{enumerate}
\hfill [JAM 2012]




\item  Find the area of the portion of the surface $z=x^{2}-y^{2}$ in $\mathbb{R}^{3}$ which lies inside the solid cylinder $x^{2}+y^{2} \leq 1$.
\hfill [JAM 2012]
\item  Let $y(x)$ be the solution of the differential equation $\frac{d^{2} y}{d x^{2}}-y=0$ such that $y(0)=2$ and $y^{\prime}(0)=2 \alpha$. Find all values of $\alpha \in[0,1)$ such that the infimum of the set $\{y(x) \mid x \in \mathbb{R}\}$ is greater than or equal to 1.
\hfill [JAM 2012]


\item 

\begin{enumerate}
  \item Assume that $y_{1}(x)=x$ and $y_{2}(x)=x^{3}$ are two linearly independent solutions of the homogeneous differential equation $x^{2} \frac{d^{2} y}{d x^{2}}-3 x \frac{d y}{d x}+3 y=0$. Using the method of variation of parameters, find a particular solution of the differential equation $x^{2} \frac{d^{2} y}{d x^{2}}-3 x \frac{d y}{d x}+3 y=x^{5}$.

\item Solve the differential equation $\frac{d y}{d x}+\frac{5 y}{6 x}=\frac{5 x^{4}}{y^{5}}$ subject to the condition $y(1)=1$.
  
\end{enumerate}
\hfill [JAM 2012]





\item

\begin{enumerate}
    
  \item Let $\vec{r}=x \hat{i}+y \hat{j}+z \hat{k}$ be the position vector field in $\mathbb{R}^{3}$ and let $f: \mathbb{R} \rightarrow \mathbb{R}$ be a differentiable function. Show that $\vec{\nabla} \times\{f(|\vec{r}|) \vec{r}\}=\overrightarrow{0}$ for $\vec{r} \neq \overrightarrow{0}$.

\item Let $W$ be the region inside the solid cylinder $x^{2}+y^{2} \leq 4$ between the plane $z=0$ and the paraboloid $z=x^{2}+y^{2}$. Let $S$ be the boundary of $W$. Using Gauss's divergence theorem, evaluate $\iint_{S} \vec{F} \cdot \hat{n} d S$, where
$$
\vec{F}=\left(x^{2}+y^{2}-4\right) \hat{i}+(3 x y) \hat{j}+\left(2 x z+z^{2}\right) \hat{k}
$$
and $\hat{n}$ is the outward unit normal vector to $S$.
\end{enumerate}
\hfill [JAM 2012]



\item 

\begin{enumerate}
     \item Let $G$ be a finite group whose order is not divisible by 3 . Show that for every $g \in G$, there exists an $h \in G$ such that $g=h^{3}$.

\item Let $A$ be the group of all rational numbers under addition, $B$ be the group of all non-zero rational numbers under multiplication and $C$ the group of all positive rational numbers under multiplication. Show that no two of the groups $A, B$ and $C$ are isomorphic.
\end{enumerate}
\hfill [JAM 2012]




\item 


\begin{enumerate}
    \item Let $I$ be an ideal of a commutative ring $R$. Define

Show that $A$ is an ideal of $R$.
$$
A=\left\{r \in R \mid r^{n} \in I \text { for some } n \in \mathbb{N}\right\} .
$$
\item Let $F$ be a field. For each $p(x) \in F[x]$ (the polynomial ring in $x$ over $F$ ) define $\varphi: F[x] \rightarrow F \times F$ by $\varphi(p(x))=(p(0), p(1))$.


\begin{enumerate}
    \item Prove that $\varphi$ is a ring homomorphism.

\item  Prove that the quotient ring $F[x] /\left(x^{2}-x\right)$ is isomorphic to the ring $F \times F$.
\end{enumerate}

\end{enumerate}
\hfill [JAM 2012]




\item 

\begin{enumerate}
     \item Let $P, D$ and $A$ be real square matrices of the same order such that $P$ is invertible, $D$ is diagonal and $D=P A P^{-1}$. If $A^{n}=0$ for some $n \in \mathbb{N}$, then show that $A=0$.

\item Let $T: V \rightarrow W$ be a linear transformation of vector spaces. Prove the following:


\begin{enumerate}
    \item  If $\left\{v_{1}, v_{2}, \ldots, v_{k}\right\}$ spans $V$, and $T$ is onto, then $\left\{T v_{1}, T v_{2}, \ldots, T v_{k}\right\}$ spans $W$.

\item  If $\left\{v_{1}, v_{2}, \ldots, v_{k}\right\}$ is linearly independent in $V$, and $T$ is one-one, then $\left\{T v_{1}, T v_{2}, \ldots, T v_{k}\right\}$ is linearly independent in $W$.

\item  If $\left\{v_{1}, v_{2}, \ldots, v_{k}\right\}$ is a basis of $V$, and $T$ is bijective, then $\left\{T v_{1}, T v_{2}, \ldots, T v_{k}\right\}$ is a basis of $W$.
\end{enumerate}



\end{enumerate}
\hfill [JAM 2012]


\item 

\begin{enumerate}
    

 \item Let $\left\{v_{1}, v_{2}, v_{3}\right\}$ be a basis of a vector space $V$ over $\mathbb{R}$. Let $T: V \rightarrow V$ be the linear transformation determined by
$$
T v_{1}=v_{1}, T v_{2}=v_{2}-v_{3} \text { and } T v_{3}=v_{2}+2 v_{3} .
$$
Find the matrix of the transformation $T$ with $\left\{v_{1}+v_{2}, v_{1}-v_{2}, v_{3}\right\}$ as a basis of both the domain and the co-domain of $T$.

\item Let $W$ be a three dimensional vector space over $\mathbb{R}$ and let $S: W \rightarrow W$ be a linear transformation. Further, assume that every non-zero vector of $W$ is an eigenvector of $S$. Prove that there exists an $\alpha \in \mathbb{R}$ such that $S=\alpha I$, where $I: W \rightarrow W$ is the identity transformation.

\end{enumerate}
\hfill [JAM 2012]

\item 

\begin{enumerate}
    
 \item Show that the function $f: \mathbb{R} \rightarrow \mathbb{R}$, defined by $f(x)=x^{2}$ for $x \in \mathbb{R}$, is not uniformly continuous.

\item For each $n \in \mathbb{N}$, let $f_{n}: \mathbb{R} \rightarrow \mathbb{R}$ be a uniformly continuous function. If the sequence $\left\{f_{n}\right\}$ converges uniformly on $\mathbb{R}$ to a function $f: \mathbb{R} \rightarrow \mathbb{R}$, then show that $f$ is uniformly continuous.

\end{enumerate}
\hfill [JAM 2012]


\item 

\begin{enumerate}
    
 \item Let $A$ be a nonempty bounded subset of $\mathbb{R}$. Show that $\{x \in \mathbb{R} \mid x \geq a$ for all $a \in A\}$ is a closed subset of $\mathbb{R}$.

\item Let $\left\{x_{n}\right\}$ be a sequence in $\mathbb{R}$ such that $\left|x_{n+1}-x_{n}\right|<\frac{1}{n^{2}}$ for all $n \in \mathbb{N}$. Show that the sequence $\left\{x_{n}\right\}$ is convergent.
\end{enumerate}
\hfill [JAM 2012]








  
\end{enumerate}

\newpage
\section{JAM 2013}
\begin{enumerate}


\item  Let $A=\left(\begin{array}{rrr}1 & 1 & 1 \\  3 & -1 & 1 \\  1 & 5 & 3\end{array}\right)$ and $V$ be the vector space of all $X \in \mathbb{R}^{3}$ such that $A X=0$. Then $\operatorname{dim}(V)$ 
is 


\begin{enumerate}
    \item $ 0$ 
\item 1 
\item 2 
\item 3
\end{enumerate}
\hfill [JAM 2013]





\item  The value of $n$ for which the divergence of the function $\vec{F}=\frac{\vec{r}}{|\vec{r}|^{n}}, \bar{r}=x \hat{i}+y \hat{j}+z \hat{k},|\vec{r}| \neq 0$, vanishes is 

\begin{enumerate} 
\item 1 
\item $-1$ 
\item 3 
\item $-3$

\end{enumerate}
\hfill [JAM 2013]


\item  Let $A$ and $B$ be subsets of $\mathbb{R}$. Which of the following is NOT necessarily true? 


\begin{enumerate}
\item $(A \cap B)^{o} \subseteq A^{o} \cap B^{o}$ 
\item $\quad A^{o} \cup B^{o} \subseteq(A \cup B)^{o}$ 
\item $\bar{A} \cup \bar{B} \subseteq \overline{A \cup B}$ 
\item $\bar{A} \cap \bar{B} \subseteq \overline{A \cap B}$

\end{enumerate}
\hfill [JAM 2013]


\item  Let $[x]$ denote the greatest integer function of $x$. The value of $\alpha$ for which the function $f(x)= \begin{cases}\frac{\sin \left[-x^{2}\right]}{\left[-x^{2}\right]}, & x \neq 0 \\ \alpha, & x=0\end{cases}$

is continuous at $x=0$ is 


\begin{enumerate}
    
\item 0 
\item $\sin (-1)$ 
\item $\sin 1$ 
\item 1
\end{enumerate}
\hfill [JAM 2013]



\item  Let the function $f(x)$ be defined by $f(x)= \begin{cases}e^{x}, & x \text { is rational } \\ e^{1-x}, & x \text { is irrational }\end{cases}$ for $x$ in $(0,1)$. Then 


\begin{enumerate}
    
\item $f$ is continuous at every point in $(0,1)$ 
\item $f$ is discontinuous at every point in $(0,1)$ 
\item $f$ is discontinuous only at one point in $(0,1)$ 
\item $f$ is continuous only at one point in $(0,1)$
\end{enumerate}
\hfill [JAM 2013]



\item  The value of the integral
$$
\iint_{D} \sqrt{x^{2}+y^{2}} d x d y, \quad D=\left\{(x, y) \in \mathbb{R}^{2}: \quad x \leq x^{2}+y^{2} \leq 2 x\right\}
$$
is 


\begin{enumerate}
   
\item 0 
\item $\frac{7}{9}$ 
\item $\frac{14}{9}$ 
\item $\frac{28}{9}$ 
\end{enumerate}
\hfill [JAM 2013]



\item  Let
$$
x_{n}=\left(1-\frac{1}{3}\right)^{2}\left(1-\frac{1}{6}\right)^{2}\left(1-\frac{1}{10}\right)^{2} \ldots\left(1-\frac{1}{\frac{n(n+1)}{2}}\right)^{2}, \quad n \geq 2 .
$$
Then $\lim _{n \rightarrow \infty} x_{n}$ is 

\begin{enumerate}
 
\item $\frac{1}{3}$ 
\item $\frac{1}{9}$ 
\item $\frac{1}{81}$ 
\item 0
\end{enumerate}
\hfill [JAM 2013]



\item  Let $p$ be a prime number. Let $G$ be the group of all $2 \times 2$ matrices over $\mathbb{Z}_{p}$ with determinant 1 under matrix multiplication. Then the order of $G$ is 

\begin{enumerate}
 \item $\quad(p-1) p(p+1)$ 
\item $\quad p^{2}(p-1)$ 
\item $p^{3}$ 
\item $p^{2}(p-1)+p$
\end{enumerate}
\hfill [JAM 2013]


 \item  Let $V$ be the vector space of all $2 \times 2$ matrices over $\mathbb{R}$. Consider the subspaces
$$
W_{1}=\left\{\left(\begin{array}{cc}
a & -a  \\
c & d
\end{array}\right): a, c, d \in \mathbb{R}\right\} \text { and } W_{2}=\left\{\left(\begin{array}{rr}
a & b  \\
-a & d
\end{array}\right): a, b, d \in \mathbb{R}\right\} \text {. }
$$
If $m=\operatorname{dim}\left(W_{1} \cap W_{2}\right)$ and $n=\operatorname{dim}\left(W_{1}+W_{2}\right)$, then the pair $(m, n)$ is 

\begin{enumerate}
 \item $(2,3)$ 
\item $(2,4)$ 
\item $(3,4)$ 
\item $(1,3)$
\end{enumerate}
\hfill [JAM 2013]




\item  Let $\wp_{n}$ be the real vector space of all polynomials of degree at most $n$. Let $D: \wp_{n} \rightarrow \wp_{n-1}$ and $T: \wp_{n} \rightarrow \wp_{n+1}$ be the linear transformations defined by

$D\left(a_{0}+a_{1} x+a_{2} x^{2}+\ldots+a_{n} x^{n}\right)=a_{1}+2 a_{2} x+\ldots+n a_{n} x^{n-1}$,

$T\left(a_{0}+a_{1} x+a_{2} x^{2}+\ldots+a_{n} x^{n}\right)=a_{0} x+a_{1} x^{2}+a_{2} x^{3}+\ldots+a_{n} x^{n+1}$,

respectively. If $A$ is the matrix representation of the transformation $D T-T D: \wp_{n} \rightarrow \wp_{n}$ with respect to the standard basis of $\wp_{n}$, then the trace of $A$ is 

\begin{enumerate}
 \item $-n$ 
\item $n$ 
\item $n+1$ 
\item $-(n+1)$
\end{enumerate}
\hfill [JAM 2013]




\item  The equation of the curve satisfying $\sin y \frac{d y}{d x}=\cos y(1-x \cos y)$ and passing through the origin is
\hfill [JAM 2013]



\item  Let $f$ be a continuously differentiable function such that $\int_{0}^{2 x^{2}} f(t) d t=e^{\cos x^{2}}$ for all $x \in(0, \infty)$. The value of $f^{\prime}(\pi)$ is
\hfill [JAM 2013]



\item  Let $u=\frac{y^{2}-x^{2}}{x^{2} y^{2}}, v=\frac{z^{2}-y^{2}}{y^{2} z^{2}}$ for $x \neq 0, y \neq 0, z \neq 0$. Let $w=f(u, v)$, where $f$ is a real valued function defined on $\mathbb{R}^{2}$ having continuous first order partial derivatives. The value of $x^{3} \frac{\partial w}{\partial x}+y^{3} \frac{\partial w}{\partial y}+z^{3} \frac{\partial w}{\partial z}$ at the point $(1,2,3)$ is
\hfill [JAM 2013]



\item  The set of points at which the function $f(x, y)=x^{4}+y^{4}-x^{2}-y^{2}+1,(x, y) \in \mathbb{R}^{2}$ attains local maximum is
\hfill [JAM 2013]



\item  Let $C$ be the boundary of the region in the first quadrant bounded by $y=1-x^{2}, x=0$ and $y=0$, oriented counter-clockwise. The value of $\oint_{C}\left(x y^{2} d x-x^{2} y d y\right)$ is
\hfill [JAM 2013]



\item  Let $f(x)=\left\{\begin{array}{cc}0, & -1  \leq x \leq 0  \\ x^{4}, & 0<x \leq 1\end{array}\right.$. If $f(x)=\sum_{k=0}^{n} \frac{f^{(k)}(0)}{k !} x^{k}+\frac{f^{(n+1)}(\xi)}{(n+1) !} x^{n+1}$ is the Taylor's formula for $f$ about $x=0$ with maximum possible value of $n$, then the value of $\xi$ for $0<x \leq 1$ is
\hfill [JAM 2013]



\item  Let $\vec{F}=2 z \hat{i}+4 x \hat{j}+5 y \hat{k}$, and let $C$ be the curve of intersection of the plane $z=x+4$ and the cylinder $x^{2}+y^{2}=4$, oriented counter-clockwise. The value of $\oint_{C} \vec{F} \cdot d \vec{r}$ is
\hfill [JAM 2013]



\item  Let $f$ and $g$ be the functions from $\mathbb{R} \backslash\{0,1\}$ to $\mathbb{R}$ defined by $f(x)=\frac{1}{x}$ and $g(x)=\frac{x-1}{x}$ for $x \in \mathbb{R} \backslash\{0,1\}$. The smallest group of functions from $\mathbb{R} \backslash\{0,1\}$ to $\mathbb{R}$ containing $f$ and $g$ under composition of functions is isomorphic to
\hfill [JAM 2013]



\item  The orthogonal trajectory of the family of curves $\frac{x^{2}}{2}+y^{2}=c$, which passes through $(1,1)$ is \item  The function to which the power series $\sum_{n=1}^{\infty}(-1)^{n+1} n x^{2 n-2}$ converges is
\hfill [JAM 2013]



\item  Let $0<a \leq 1, s_{1}=\frac{a}{2}$ and for $n \in \mathbb{N}$, let $s_{n+1}=\frac{1}{2}\left(s_{n}^{2}+a\right)$. Show that the sequence $\left\{s_{n}\right\}$ is convergent, and find its limit.
\hfill [JAM 2013]



\item  Evaluate

by changing the order of integration.
$$
\int_{1 / 4}^{1} \int_{\sqrt{x-x^{2}}}^{\sqrt{x}} \frac{x^{2}-y^{2}}{x^{2}} d y d x
$$
\hfill [JAM 2013]


\item  Find the general solution of the differential equation
$$
x^{2} \frac{d^{3} y}{d x^{3}}+x \frac{d^{2} y}{d x^{2}}-6 \frac{d y}{d x}+6 \frac{y}{x}=\frac{x \ln x+1}{x^{2}}, x>0 .
$$
\hfill [JAM 2013]


\item  Let $S_{1}$ be the hemisphere $x^{2}+y^{2}+z^{2}=1, z>0$ and $S_{2}$ be the closed disc $x^{2}+y^{2} \leq 1$ in the $x y$ plane. Using Gauss' divergence theorem, evaluate $\iint_{S} \vec{F} \cdot d \bar{S}$, where $\vec{F}=z^{2} x \hat{i}+\left(\frac{y^{3}}{3}+\tan z\right) \hat{j}+\left(x^{2} z+y^{2}\right) \hat{k}$ and $S=S_{1} \cup S_{2}$.

Also evaluate $\iint_{S_{1}} \vec{F} \cdot d \vec{S}$.

\hfill [JAM 2013]


\item  Let
$$
f(x, y)= \begin{cases}\frac{2\left(x^{3}+y^{3}\right)}{x^{2}+2 y}, & (x, y) \neq(0,0) \\ 0, & (x, y)=(0,0) .\end{cases}
$$
Show that the first order partial derivatives of $f$ with respect to $x$ and $y$ exist at $(0,0)$. Also show that $f$ is not continuous at $(0,0)$.
\hfill [JAM 2013]



\item  Let $A$ be an $n \times n$ diagonal matrix with characteristic polynomial $(x-a)^{p}(x-b)^{q}$, where $a$ and $b$ are distinct real numbers. Let $V$ be the real vector space of all $n \times n$ matrices $B$ such that $A B=B A$. Determine the dimension of $V$.
\hfill [JAM 2013]



\item  Let $A$ be an $n \times n$ real symmetric matrix with $n$ distinct eigenvalues. Prove that there exists an orthogonal matrix $P$ such that $A P=P D$, where $D$ is a real diagonal matrix.
\hfill [JAM 2013]



\item  Let $K$ be a compact subset of $\mathbb{R}$ with nonempty interior. Prove that $K$ is of the form $[a, b]$ or of the form $[a, b] \backslash \bigcup I_{n}$, where $\left\{I_{n}\right\}$ is a countable disjoint family of open intervals with end points in $K$.
\hfill [JAM 2013]



\item  Let $f:[a, b] \rightarrow \mathbb{R}$ be a continuous function such that $f$ is differentiable in $(a, c)$ and $(c, b)$, $a<c<b$. If $\lim _{x \rightarrow c} f^{\prime}(x)$ exists, then prove that $f$ is differentiable at $c$ and $f^{\prime}(c)=\lim _{x \rightarrow c} f^{\prime}(x)$.
\hfill [JAM 2013]



\item  Let $G$ be a finite group, and let $\varphi$ be an automorphism of $G$ such that $\varphi(x)=x$ if and only if $x=e$, where $e$ is the identity element in $G$. Prove that every $g \in G$ can be represented as $g=x^{-1} \varphi(x)$ for some $x \in G$. Moreover, if $\varphi(\varphi(x))=x$ for every $x \in G$, then show that $G$ is abelian.
\hfill [JAM 2013]




    
\end{enumerate}
\newpage
\section{JAM 2014}

\begin{enumerate}



\item  Let $f(x)=\left|x^{2}-25\right|$ for all $x \in \mathbb{R}$. The total number of points of $\mathbb{R}$ at which $f$ attains a local extremum (minimum or maximum) is 

\begin{enumerate}
 \item 1 
\item 2 
\item 3 
\item 4

\end{enumerate}
\hfill [JAM 2014]


\item  The coefficient of $(x-1)^{2}$ in the Taylor series expansion of $f(x)=x e^{x}(x \in \mathbb{R})$ about the point $x=1$ is 

\begin{enumerate}
 \item $\frac{e}{2}$ 
\item $2 e$ 
\item $\frac{3 e}{2}$ 
\item $3 e$
\end{enumerate}
\hfill [JAM 2014]



\item  Let $f(x, y)=\sum_{k=1}^{10}\left(x^{2}-y^{2}\right)^{k}$ for all $(x, y) \in \mathbb{R}^{2}$. Then for all $(x, y) \in \mathbb{R}^{2}$, 

\begin{enumerate}
 \item $x \frac{\partial f}{\partial x}(x, y)-y \frac{\partial f}{\partial y}(x, y)=0$ 
\item $x \frac{\partial f}{\partial x}(x, y)+y \frac{\partial f}{\partial y}(x, y)=0$ 
\item $y \frac{\partial f}{\partial x}(x, y)-x \frac{\partial f}{\partial y}(x, y)=0$ 
\item $y \frac{\partial f}{\partial x}(x, y)+x \frac{\partial f}{\partial y}(x, y)=0$
\end{enumerate}
\hfill [JAM 2014]



\item  For $a, b, c \in \mathbb{R}$, if the differential equation
$$
\left(a x^{2}+b x y+y^{2}\right) d x+\left(2 x^{2}+c x y+y^{2}\right) d y=0
$$
is exact, then 

\begin{enumerate}
 \item $b=2, c=2 a$ 
\item $b=4, c=2$ 
\item $b=2, c=4$ 
\item $b=2, a=2 c$
\end{enumerate}
\hfill [JAM 2014]



\item  If $f(x, y, z)=x^{2} y+y^{2} z+z^{2} x$ for all $(x, y, z) \in \mathbb{R}^{3}$ and $\nabla=\frac{\partial}{\partial x} \mathbf{i}+\frac{\partial}{\partial y} \mathbf{j}+\frac{\partial}{\partial z} \mathbf{k}$, then the value of $\nabla \cdot(\nabla \times \nabla f)+\nabla \cdot(\nabla f)$ at $(1,1,1)$ is 

\begin{enumerate}
 \item 0 
\item 3 
\item 6 
\item 9
\end{enumerate}
\hfill [JAM 2014]



\item  The radius of convergence of the power series $\sum_{n=0}^{\infty} 2^{2 n} x^{n^{2}}$ is 

\begin{enumerate}
 \item $\frac{1}{4}$ 
\item 1 
\item 2 
\item 4
\end{enumerate}
\hfill [JAM 2014]



\item  Let $G$ be a group of order 17. The total number of non-isomorphic subgroups of $G$ is 

\begin{enumerate}
 \item 1 
\item 2 
\item 3 
\item 17
\end{enumerate}
\hfill [JAM 2014]



\item  Which one of the following is a subspace of the vector space $\mathbb{R}^{3}$ ? 

\begin{enumerate}
 
\item $\left\{(x, y, z) \in \mathbb{R}^{3}: x+2 y=0,2 x+3 z=0\right\}$ 
\item $\left\{(x, y, z) \in \mathbb{R}^{3}: 2 x+3 y+4 z-3=0, z=0\right\}$ 
\item $\left\{(x, y, z) \in \mathbb{R}^{3}: x \geq 0, y \geq 0\right\}$ 
\item $\left\{(x, y, z) \in \mathbb{R}^{3}: x-1=0, y=0\right\}$

\end{enumerate}
\hfill [JAM 2014]



\item  Let $T: \mathbb{R}^{3} \rightarrow \mathbb{R}^{3}$ be the linear transformation defined by $T(x, y, z)=(x+y, y+z, z+x)$ for all $(x, y, z) \in \mathbb{R}^{3}$. Then 

\begin{enumerate}
 
\item $\operatorname{rank}(T)=0, \operatorname{nullity}(T)=3$ 
\item $\operatorname{rank}(T)=2, \operatorname{nullity}(T)=1$ 
\item $\operatorname{rank}(T)=1$, $\operatorname{nullity}(T)=2$ 
\item $\operatorname{rank}(T)=3, \operatorname{nullity}(T)=0$

\end{enumerate}
\hfill [JAM 2014]



\item  Let $f: \mathbb{R} \rightarrow \mathbb{R}$ be a continuous function satisfying $x+\int_{0}^{x} f(t) d t=e^{x}-1$ for all $x \in \mathbb{R}$. Then the set $\{x \in \mathbb{R}: 1 \leq f(x) \leq 2\}$ is the interval 

\begin{enumerate}
 \item $[\log 2, \log 3]$ 
\item $[2 \log 2,3 \log 3]$ 
\item $\left[e-1, e^{2}-1\right]$ 
\item $\left[0, e^{2}\right]$
\end{enumerate}
\hfill [JAM 2014]



 \item  The system of linear equations
$$
\begin{aligned}
x-y+2 z &=b_{1} \\
x+2 y-z &=b_{2} \\
2 y-2 z &=b_{3}
\end{aligned}
$$
is inconsistent when $\left(b_{1}, b_{2}, b_{3}\right)$ equals 

\begin{enumerate}
 \item $(2,2,0)$ 
\item $(0,3,2)$ 
\item $\quad(2,2,1)$ 
\item $(2,-1,-2)$
\end{enumerate}
\hfill [JAM 2014]



\item  Let $A=\left[\begin{array}{ccc}a & -1 & 4 \\ 0 & b & 7 \\ 0 & 0 & 3\end{array}\right]$ be a matrix with real entries. If the sum and the product of all the eigenvalues of $A$ are 10 and 30 respectively, then $a^{2}+b^{2}$ equals 


\begin{enumerate}
 \item 29 
\item 40 
\item 58 
\item 65
\end{enumerate}
\hfill [JAM 2014]



\item  Consider the subspace $W=\left\{\left(x_{1}, x_{2}, \ldots, x_{10}\right) \in \mathbb{R}^{10}: x_{n}=x_{n-1}+x_{n-2}\right.$ for $\left.3 \leq n \leq 10\right\}$ of the vector space $\mathbb{R}^{10}$. The dimension of $W$ is 

\begin{enumerate}
 
\item 2 
\item 3 
\item 9 
\item 10
\end{enumerate}
\hfill [JAM 2014]


\item  Let $y_{1}(x)$ and $y_{2}(x)$ be two linearly independent solutions of the differential equation
$$
x^{2} y^{\prime \prime}(x)-2 x y^{\prime}(x)-4 y(x)=0 \text { for } x \in[1,10] \text {. }
$$
Consider the Wronskian $W(x)=y_{1}(x) y_{2}^{\prime}(x)-y_{2}(x) y_{1}^{\prime}(x)$. If $W(1)=1$, then $W(3)-W(2)$ equals 


\begin{enumerate}
 
\item 1 
\item 2 
\item 3 
\item 5
\end{enumerate}
\hfill [JAM 2014]


\item  The equation of the curve passing through the point $\left(\frac{\pi}{2}, 1\right)$ and having slope $\frac{\sin (x)}{x^{2}}-\frac{2 y}{x}$ at each point $(x, y)$ with $x \neq 0$ is 


\begin{enumerate}
 \item $-x^{2} y+\cos (x)=\frac{-\pi^{2}}{4}$ 
\item $\quad x^{2} y+\cos (x)=\frac{\pi^{2}}{4}$ 
\item $x^{2} y-\sin (x)=\frac{\pi^{2}}{4}-1$ 
\item $x^{2} y+\sin (x)=\frac{\pi^{2}}{4}+1$
\end{enumerate}
\hfill [JAM 2014]



\item  The value of $\alpha \in \mathbb{R}$ for which the curves $x^{2}+\alpha y^{2}=1$ and $y=x^{2}$ intersect orthogonally is 


\begin{enumerate}
 \item $-2$ 
\item $\frac{-1}{2}$ 
\item $\frac{1}{2}$ 
\item 2
\end{enumerate}
\hfill [JAM 2014]



\item  Let $x_{n}=2^{2 n}\left(1-\cos \left(\frac{1}{2^{n}}\right)\right)$ for all $n \in \mathbb{N}$. Then the sequence $\left\{x_{n}\right\}$ 

\begin{enumerate}
 \item does NOT converge 
\item converges to 0 
\item converges to $\frac{1}{2}$ 
\item converges to $\frac{1}{4}$
\end{enumerate}
\hfill [JAM 2014]



\item  Let $\left\{x_{n}\right\}$ be a sequence of real numbers such that $\lim _{n \rightarrow \infty}\left(x_{n+1}-x_{n}\right)=c$, where $c$ is a positive real number. Then the sequence $\left\{\frac{x_{n}}{n}\right\}$ 

\begin{enumerate}
 \item is NOT bounded 
\item is bounded but NOT convergent 
\item converges to $c$ 
\item converges to 0
\end{enumerate}
\hfill [JAM 2014]

 \item  Let $\sum_{n=1}^{\infty} a_{n}$ and $\sum_{n=1}^{\infty} b_{n}$ be two series, where $a_{n}=\frac{(-1)^{n} n}{2^{n}}, b_{n}=\frac{(-1)^{n}}{\log (n+1)}$ for all $n \in \mathbb{N}$. Then
 

\begin{enumerate}
 \item both $\sum_{n=1}^{\infty} a_{n}$ and $\sum_{n=1}^{\infty} b_{n}$ are absolutely convergent 
\item $\sum_{n=1}^{\infty} a_{n}$ is absolutely convergent but $\sum_{n=1}^{\infty} b_{n}$ is conditionally convergent 
\item $\sum_{n=1}^{\infty} a_{n}$ is conditionally convergent but $\sum_{n=1}^{\infty} b_{n}$ is absolutely convergent 
\item both $\sum_{n=1}^{\infty} a_{n}$ and $\sum_{n=1}^{\infty} b_{n}$ are conditionally convergent
\end{enumerate}
\hfill [JAM 2014]



\item  The set $\left\{\frac{x^{2}}{1+x^{2}}: x \in \mathbb{R}\right\}$ is

\begin{enumerate}
 \item connected but NOT compact in $\mathbb{R}$ 
\item compact and connected in $\mathbb{R}$ 
\item compact but NOT connected in $\mathbb{R}$ 
\item neither compact nor connected in $\mathbb{R}$
\end{enumerate}
\hfill [JAM 2014]
 


\item  The set of all limit points of the set $\left\{\frac{2}{x+1}: x \in(-1,1)\right\}$ in $\mathbb{R}$ is 

\begin{enumerate}
 
\item $[1, \infty)$ 
\item $(1, \infty)$ 
\item $[-1,1]$ 
\item $[-1, \infty)$
\end{enumerate}
\hfill [JAM 2014]


\item  Let $S=[0,1] \cup[2,3)$ and let $f: S \rightarrow \mathbb{R}$ be defined by $f(x)=\left\{\begin{array}{lll}2 x & \text { if } & x \in[0,1] \text {, } \\ 8-2 x & \text { if } & x \in[2,3) .\end{array}\right.$ If $T=\{f(x): x \in S\}$, then the inverse function $f^{-1}: T \rightarrow S$ 

\begin{enumerate}
 \item does NOT exist 
\item exists and is continuous 
\item exists and is NOT continuous 
\item exists and is monotonic
\end{enumerate}
\hfill [JAM 2014]



\item  Let $f(x)=x^{3}+x$ and $g(x)=x^{3}-x$ for all $x \in \mathbb{R}$. If $f^{-1}$ denotes the inverse function of $f$, then the derivative of the composite function $g \circ f^{-1}$ at the point 2 is 

\begin{enumerate}
 
\item $\frac{2}{13}$ 
\item $\frac{1}{2}$ 
\item $\frac{11}{13}$ 
\item $\frac{11}{4}$
\end{enumerate}
\hfill [JAM 2014]


\item  For all $(x, y) \in \mathbb{R}^{2}$, let $f(x, y)= \begin{cases}x & \text { if } y=0, \\ x-y^{3} \sin (1 / y) & \text { if } y \neq 0 .\end{cases}$

Then at the point $(0,0)$, 


\begin{enumerate}
 \item $f$ is NOT continuous 
\item $f$ is continuous but NOT differentiable 
\item $\frac{\partial f}{\partial x}$ exists but $\frac{\partial f}{\partial y}$ does NOT exist 
\item $f$ is differentiable
\end{enumerate}
\hfill [JAM 2014]



\item  For all $(x, y) \in \mathbb{R}^{2}$, let $f(x, y)= \begin{cases}\frac{x}{|x|} \sqrt{x^{2}+y^{2}} & \text { if } x \neq 0, \\ 0 & \text { if } x=0 .\end{cases}$

Then $\frac{\partial f}{\partial x}(0,0)+\frac{\partial f}{\partial y}(0,0)$ equals 


\begin{enumerate}
 
\item $-1$ 
\item 0 
\item 1 
\item 2
\end{enumerate}
\hfill [JAM 2014]





 \item  Let $f: \mathbb{R} \rightarrow \mathbb{R}$ be a function with continuous derivative such that $f(\sqrt{2})=2$ and $f(x)=\lim _{t \rightarrow 0} \frac{1}{2 t} \int_{x-t}^{x+t} s f^{\prime}(s) d s$ for all $x \in \mathbb{R}$. Then $f(3)$ equals 
 

\begin{enumerate}
 \item $\sqrt{3}$ 
\item $\quad 3 \sqrt{2}$ 
\item $3 \sqrt{3}$ 
\item 9

\end{enumerate}
\hfill [JAM 2014]


\item  The value of $\int_{x=0}^{1} \int_{y=0}^{x^{2}} \int_{z=0}^{y}(y+2 z) d z d y d x$ is 


\begin{enumerate}
 
\item $\frac{1}{53}$ 
\item $\frac{2}{21}$ 
\item $\frac{1}{6}$ 
\item $\frac{5}{3}$
\end{enumerate}
\hfill [JAM 2014]


\item  If $C$ is a smooth curve in $\mathbb{R}^{3}$ from $(-1,0,1)$ to $(1,1,-1)$, then the value of $\int_{C}\left(2 x y+z^{2}\right) d x+\left(x^{2}+z\right) d y+(y+2 x z) d z$ is 

\begin{enumerate}
 \item 0 
\item 1 
\item 2 
\item 3

\end{enumerate}
\hfill [JAM 2014]


\item  Let $C$ be the boundary of the region $R=\left\{(x, y) \in \mathbb{R}^{2}:-1 \leq y \leq 1,0 \leq x \leq 1-y^{2}\right\}$ oriented in the counterclockwise direction. Then the value of $\oint_{C} y d x+2 x d y$ is 

\begin{enumerate}
 \item $\frac{-4}{3}$ 
\item $\frac{-2}{3}$ 
\item $\frac{2}{3}$ 
\item $\frac{4}{3}$
\end{enumerate}
\hfill [JAM 2014]



\item  Let $G$ be a cyclic group of order 24 . The total number of group isomorphisms of $G$ onto itself is 

\begin{enumerate}
 \item 7 
\item 8 
\item 17 
\item 24
\end{enumerate}
\hfill [JAM 2014]



\item  Let $S_{n}$ be the group of all permutations on the set $\{1,2, \ldots, n\}$ under the composition of mappings. For $n>2$, if $H$ is the smallest subgroup of $S_{n}$ containing the transposition $(1,2)$ and the cycle $(1,2, \ldots, n)$, then 

\begin{enumerate}
 
\item $H=S_{n}$ 
\item $H$ is abelian 
\item the index of $H$ in $S_{n}$ is 2 
\item $\mathrm{H}$ is cyclic
\end{enumerate}
\hfill [JAM 2014]


\item  Let $S$ be the oriented surface $x^{2}+y^{2}+z^{2}=1$ with the unit normal $\mathbf{n}$ pointing outward. For the vector field $\mathbf{F}(x, y, z)=x \mathbf{i}+y \mathbf{j}+z \mathbf{k}$, the value of $\iint_{S} \mathbf{F} \cdot \mathbf{n} d S$ is 

\begin{enumerate}
 \item $\frac{\pi}{3}$ 
\item $2 \pi$ 
\item $\frac{4 \pi}{3}$ 
\item $4 \pi$
\end{enumerate}
\hfill [JAM 2014]



\item  Let $f:(0, \infty) \rightarrow \mathbb{R}$ be a differentiable function such that $f^{\prime}\left(x^{2}\right)=1-x^{3}$ for all $x>0$ and $f(1)=0$. Then $f(4)$ equals 

\begin{enumerate}
 \item $\frac{-47}{5}$ 
\item $\frac{-47}{10}$ 
\item $\frac{-16}{5}$ 
\item $\frac{-8}{5}$
\end{enumerate}
\hfill [JAM 2014]

 \item  Which one of the following conditions on a group $G$ implies that $G$ is abelian? 

\begin{enumerate}
 \item The order of $G$ is $p^{3}$ for some prime $p$ 
\item Every proper subgroup of $G$ is cyclic 
\item Every subgroup of $G$ is normal in $G$ 
\item The function $f: G \rightarrow G$, defined by $f(x)=x^{-1}$ for all $x \in G$, is a homomorphism
\end{enumerate}
\hfill [JAM 2014]



\item  Let $S=\left\{x \in \mathbb{R}: x^{6}-x^{5} \leq 100\right\}$ and $T=\left\{x^{2}-2 x: x \in(0, \infty)\right\}$. The set $S \cap T$ is 

\begin{enumerate}
 \item closed and bounded in $\mathbb{R}$ 
\item closed but NOT bounded in $\mathbb{R}$ 
\item bounded but NOT closed in $\mathbb{R}$ 
\item neither closed nor bounded in $\mathbb{R}$
\end{enumerate}
\hfill [JAM 2014]



\item  Find all the critical points of the function $f: \mathbb{R}^{2} \rightarrow \mathbb{R}$ defined by $f(x, y)=x^{3}+x y+y^{3}$ for all $(x, y) \in \mathbb{R}^{2}$. Also, examine whether the function $f$ attains a local maximum or a local minimum at each of these critical points.
\hfill [JAM 2014]




\item  Given that there is a common solution to the following equations:
$$
\begin{aligned}
&\mathbf{P}: y^{\prime}+2 y=e^{x} y^{2}, \quad y(0)=1,  
&\mathbf{Q}: y^{\prime \prime}-2 y^{\prime}+\alpha y=0
\end{aligned}
$$
find the value of $\alpha$ and hence find the general solution of $\mathbf{Q}$.
\hfill [JAM 2014]




\item  Let $f: \mathbb{R} \rightarrow \mathbb{R}$ be a twice differentiable function such that $f\left(\frac{1}{2^{n}}\right)=0$ for all $n \in \mathbb{N}$. Show that $f^{\prime}(0)=0=f^{\prime \prime}(0)$.
\hfill [JAM 2014]




\item  Let $A$ be an $n \times n$ matrix with real entries such that $A^{2}=A$. If I denotes the $n \times n$ identity matrix, then show that $\operatorname{rank}(A-I)=\operatorname{nullity}(A)$.
\hfill [JAM 2014]




\item  Evaluate $\iint_{S} \frac{x y}{\sqrt{1+2 x^{2}}} d S$, where the surface $S=\left\{\left(x, y, x^{2}+y\right) \in \mathbb{R}^{3}: 0 \leq x \leq y, x+y \leq 1\right\}$.
\hfill [JAM 2014]




\item  Let $f:(0,1) \rightarrow \mathbb{R}$ be a differentiable function such that $\left|f^{\prime}(x)\right| \leq 5$ for all $x \in(0,1)$. Show that the sequence $\left\{f\left(\frac{1}{n+1}\right)\right\}$ converges in $\mathbb{R}$.

\hfill [JAM 2014]



\item  Let $H$ be a subgroup of the group $(\mathbb{R},+)$ such that $H \cap[-1,1]$ is a finite set containing a nonzero element. Show that $H$ is cyclic.

\hfill [JAM 2014]



\item  If $K$ is a nonempty closed subset of $\mathbb{R}$, then show that the set $\{x+y: x \in K, y \in[1,2]\}$ is closed in $\mathbb{R}$.

\hfill [JAM 2014]





    
\end{enumerate}
\newpage
\section{JAM 2015}
\begin{enumerate}


 
\item  Suppose $N$ is a normal subgroup of a group $G$. Which one of the following is true? 
\begin{enumerate}
   \item If $G$ is an infinite group then $G / \mathrm{N}$ is an infinite group 
\item If $G$ is a nonabelian group then $G / N$ is a nonabelian group 
\item If $G$ is a cyclic group then $G / N$ is an abelian group 
\item If $G$ is an abelian group then $G / N$ is a cyclic group
 
\end{enumerate}
\hfill [JAM 2015]




\item  Let $y(x)=u(x) \sin x+v(x) \cos x$ be a solution of the differential equation $y^{\prime \prime}+y=\sec x$. Then $u(x)$ is 

\begin{enumerate}
    
\item $\ln |\cos x|+C$ 
\item $x+C$ 
\item $-x+C$ 
\item $\ln |\sec x|+C$
\end{enumerate}
\hfill [JAM 2015]



\item  Let $a, b, c, d$ be distinct non-zero real numbers with $a+b=c+d$. Then an eigenvalue of the matrix $\left[\begin{array}{ccc}a & b & 1 \\ c & d & 1 \\ 1 & -1 & 0\end{array}\right]$ is 
\begin{enumerate}
   \item $a+c$ 
\item $a+b$ 
\item $a-b$ 
\item $b-d$ 
\end{enumerate}
\hfill [JAM 2015]





\item  Let $S$ be a nonempty subset of $\mathbb{R}$. If $S$ is a finite union of disjoint bounded intervals, then which one of the following is true? 
\begin{enumerate}
    
\item If $S$ is not compact, then sup $S \notin S$ and inf $S \notin S$ 
\item Even if $\sup S \in S$ and inf $S \in S, S$ need not be compact 
\item If $\sup S \in S$ and inf $S \in S$, then $S$ is compact 
\item Even if $S$ is compact, it is not necessary that sup $S \in S$ and inf $S \in S$
\end{enumerate}
\hfill [JAM 2015]




\item  Let $\left\{x_{n}\right\}$ be a convergent sequence of real numbers. If $x_{1}>\pi+\sqrt{2}$ and $x_{n+1}=\pi+\sqrt{x_{n}-\pi}$ for $n \geq 1$, then which one of the following is the limit of this sequence? 
\begin{enumerate}
    \item $\pi+1$ 
\item $\pi+\sqrt{2}$ 
\item $\pi$ 
\item $\pi+\sqrt{\pi}$
\end{enumerate}
\hfill [JAM 2015]





\item  The volume of the portion of the solid cylinder $x^{2}+y^{2} \leq 2$ bounded above by the surface $z=x^{2}+y^{2}$ and bounded below by the $x y$-plane is 
\begin{enumerate}
    \item $\pi$ 
\item $2 \pi$ 
\item $3 \pi$ 
\item $4 \pi$
\end{enumerate}
\hfill [JAM 2015]





\item  Let $f: \mathbb{R} \rightarrow \mathbb{R}$ be a differentiable function with $f(0)=0$. If for all $x \in \mathbb{R}, 1<f^{\prime}(x)<2$, then which one of the following statements is true on $(0, \infty)$ ? 
\begin{enumerate}
    \item $f$ is unbounded 
\item $f$ has at least one zero 
\item $f$ is increasing and bounded 
\item $f$ is periodic
\end{enumerate}
\hfill [JAM 2015]





\item  If an integral curve of the differential equation $(y-x) \frac{d y}{d x}=1$ passes through $(0,0)$ and $(\alpha, 1)$, then $\alpha$ is equal to 
\begin{enumerate}
    \item $2-e^{-1}$ 
\item $1-e^{-1}$ 
\item $e^{-1}$ 
\item $1+e$
\end{enumerate}
\hfill [JAM 2015]





\item  An integrating factor of the differential equation
$$
\frac{d y}{d x}=\frac{2 x y^{2}+y}{x-2 y^{3}}
$$
\begin{enumerate}
    
\item $\frac{1}{y}$ 
\item $\frac{1}{y^{2}}$ 
\item $y$ 
\item $y^{2}$
\end{enumerate}
\hfill [JAM 2015]




\item  Let $A$ be a nonempty subset of $\mathbb{R}$. Let $I(A)$ denote the set of interior points of $A$. Then $I(A)$ can be 
\begin{enumerate}
   \item empty 
\item singleton 
\item a finite set containing more than one element 
\item countable but not finite 
\end{enumerate}
\hfill [JAM 2015]



 \item  Let $S_{3}$ be the group of permutations of three distinct symbols. The direct sum $S_{3} \oplus S_{3}$ has an 
\begin{enumerate}
    \item 4 
\item 6 
\item 9 
\item 18
\end{enumerate}
\hfill [JAM 2015]





\item  The orthogonal trajectories of the family of curves $y=C_{1} x^{3}$ are 
\begin{enumerate}
    
\item $2 x^{2}+3 y^{2}=C_{2}$ 
\item $3 x^{2}+2 y^{2}=C_{2}$ 
\item $3 x^{2}+y^{2}=C_{2}$ 
\item $x^{2}+3 y^{2}=c_{2}$

\end{enumerate}
\hfill [JAM 2015]



\item  Let $G$ be a nonabelian group. Let $\alpha \in G$ have order 4 and let $\beta \in G$ have order 3 . Then the order of the element $\alpha \beta$ in $G$ 
\begin{enumerate}
    
\item is 6 
\item is of the form $12 k$ for $k \geq 2$ 
\item is 12 
\item need not be finite

\end{enumerate}
\hfill [JAM 2015]



\item  Let $S$ be the bounded surface of the cylinder $x^{2}+y^{2}=1$ cut by the planes $z=0$ and $z=1+x$. Then the value of the surface integral $\iint_{S} 3 z^{2} d \sigma$ is equal to 
\begin{enumerate}
    \item $\int_{0}^{2 \pi}(1+\cos \theta)^{3} d \theta$ 
\item $\int_{0}^{2 \pi} \sin \theta \cos \theta(1+\cos \theta)^{2} d \theta$ 
\item $\int_{0}^{2 \pi}(1+2 \cos \theta)^{3} d \theta$ 
\item $\int_{0}^{2 \pi} \sin \theta \cos \theta(1+2 \cos \theta)^{2} d \theta$
\end{enumerate}
\hfill [JAM 2015]





\item  Suppose that the dependent variables $z$ and $w$ are functions of the independent variables $x$ and $y$, defined by the equations $f(x, y, z, w)=0$ and $g(x, y, z, w)=0$, where $f_{z} g_{w}-f_{w} g_{z}=1$. Which one of the following is correct? 
\begin{enumerate}
    
\item $z_{x}=f_{w} g_{x}-f_{x} g_{w}$ 
\item $z_{x}=f_{x} g_{w}-f_{w} g_{x}$ 
\item $z_{x}=f_{z} g_{x}-f_{x} g_{z}$ 
\item $z_{x}=f_{z} g_{w}-f_{z} g_{x}$
\end{enumerate}
\hfill [JAM 2015]




\item 

Let $A=\left[\begin{array}{cc}0 & 1-i \\ -1-i & i\end{array}\right]$ and $B=A^{T} \bar{A}$. Then 
\begin{enumerate}
    \item an eigenvalue of $B$ is purely imaginary 
\item an eigenvalue of $A$ is zero 
\item all eigenvalues of $B$ are real 
\item $A$ has a non-zero real eigenvalue
\end{enumerate}
\hfill [JAM 2015]





\item  The limit

is equal to
$$
\lim _{x \rightarrow 0+} \frac{1}{\sin ^{2} x} \int_{\frac{x}{2}}^{x} \sin ^{-1} t d t
$$
\begin{enumerate}
    \item 0 
\item $\frac{1}{8}$ 
\item $\frac{1}{4}$ 
\item $\frac{3}{8}$
\end{enumerate}
\hfill [JAM 2015]





\item  Let $P_{2}(\mathbb{R})$ be the vector space of polynomials in $x$ of degree at most 2 with real coefficients. Let $M_{2}(\mathbb{R})$ be the vector space of $2 \times 2$ real matrices. If a linear transformation $T: P_{2}(\mathbb{R}) \rightarrow M_{2}(\mathbb{R})$ is defined as

then
$$
T(f)=\left[\begin{array}{cc}
f(0)-f(2) & 0 \\
0 & f(1)
\end{array}\right]
$$

\begin{enumerate}
    \item $T$ is one-one but not onto 
\item $T$ is onto but not one-one 
\item Range $(T)=\operatorname{span}\left\{\left[\begin{array}{ll}0 & 0 \\ 0 & 1\end{array}\right],\left[\begin{array}{cc}-2 & 0 \\ 0 & 1\end{array}\right]\right\}$ 
\item Null $(T)=\operatorname{span}\left\{x^{2}-2 x, 1-x\right\}$ 

\end{enumerate}
\hfill [JAM 2015]





\item Let $B_{1}=\{(1,2),(2,-1)\}$ and $B_{2}=\{(1,0),(0,1)\}$ be ordered bases of $\mathbb{R}^{2}$. If $T: \mathbb{R}^{2} \rightarrow \mathbb{R}^{2}$ is a linear transformation such that $[T]_{B_{1}, B_{2}}$, the matrix of $T$ with respect to $B_{1}$ and $B_{2}$, is $\left[\begin{array}{cc}4 & 3 \\ 2 & -4\end{array}\right]$, then $T(5,5)$ is equal to 

\begin{enumerate}
    
\item $(-9,8)$ 
\item $(9,8)$ 
\item $(-15,-2)$ 
\item $(15,2)$
\end{enumerate}
\hfill [JAM 2015]


\item 

Let $S=\bigcap_{n=1}^{\infty}\left(\left[0, \frac{1}{2 n+1}\right] \cup\left[\frac{1}{2 n}, 1\right]\right)$. Which one of the following statements is FALSE? 

\begin{enumerate}
    \item There exist sequences $\left\{a_{n}\right\}$ and $\left\{b_{n}\right\}$ in $[0,1]$ such that $S=[0,1] \backslash \cup_{n=1}^{\infty}\left(a_{n}, b_{n}\right)$ 
\item $[0,1] \backslash S$ is an open set 
\item If $A$ is an infinite subset of $S$, then $A$ has a limit point 
\item There exists an infinite subset of $S$ having no limit points

\end{enumerate}
\hfill [JAM 2015]




\item  Let $f: \mathbb{R} \rightarrow \mathbb{R}$ be a strictly increasing continuous function. If $\left\{a_{n}\right\}$ is a sequence in $[0,1]$, then the sequence $\left\{f\left(a_{n}\right)\right\}$ is 
\begin{enumerate}
    \item increasing 
\item bounded 
\item convergent 
\item not necessarily bounded
\end{enumerate}
\hfill [JAM 2015]





\item 

Which one of the following statements is true for the series $\sum_{n=1}^{\infty}(-1)^{n} \frac{(2 n) !}{n^{2 n}}$ ? 
\begin{enumerate}
    \item The series converges conditionally but not absolutely 
\item The series converges absolutely 
\item The sequence of partial sums of the series is bounded but not convergent 
\item The sequence of partial sums of the series is unbounded

\end{enumerate}
\hfill [JAM 2015]




\item  

The sequence $\left\{\cos \left(\frac{1}{2} \tan ^{-1}\left(-\frac{n}{2}\right)^{n}\right)\right\}$ is 

\begin{enumerate}
    \item monotone and convergent 
\item monotone but not convergent 
\item convergent but not monotone 
\item neither monotone nor convergent
\end{enumerate}
\hfill [JAM 2015]





\item  

If $y(t)$ is a solution of the differential equation $y^{\prime \prime}+4 y=2 e^{t}$, then
$$
\lim _{t \rightarrow \infty} e^{-t} y(t)
$$
is equal to

\begin{enumerate}
    
\item $\frac{2}{3}$ 
\item $\frac{2}{5}$ 
\item $\frac{2}{7}$ 
\item $\frac{2}{9}$
\end{enumerate}
\hfill [JAM 2015]




\item  

For what real values of $x$ and $y$, does the integral $\int_{x}^{y}\left(6-t-t^{2}\right) d t$ attain its maximum? 

\begin{enumerate}
    
\item $x=-3, y=2$ 
\item $x=-2, y=2$ 
\item $x=2, y=3$ 
\item $x=-3, y=4$
\end{enumerate}
\hfill [JAM 2015]




\item  The area of the planar region bounded by the curves $x=6 y^{2}-2$ and $x=2 y^{2}$ is 
\begin{enumerate}
   
\item $\frac{\sqrt{2}}{3}$ 
\item $\frac{2 \sqrt{2}}{3}$ 
\item $\frac{4 \sqrt{2}}{3}$ 
\item $\sqrt{2}$  
\end{enumerate}
\hfill [JAM 2015]

\item 
For $n \geq 2$, let $f_{n}: \mathbb{R} \rightarrow \mathbb{R}$ be given by $f_{n}(x)=x^{n} \sin x$. Then at $x=0, f_{n}$ has a 
\begin{enumerate}
    \item local maximum if $n$ is even 
\item local maximum if $n$ is odd 
\item local minimum if $n$ is even 
\item local minimum if $n$ is odd

\end{enumerate}
\hfill [JAM 2015]




\item  

For $m, n \in \mathbb{N}$, define $f_{m, n}(x)=\left\{\begin{array}{c}x^{m} \sin \left(\frac{1}{x^{n}}\right. \\ 0\end{array} \quad x=0\right.$

Then at $x=0, f_{m, n}$ is 

\begin{enumerate}
    \item differentiable for each pair $m, n$ with $m>n$ 
\item differentiable for each pair $m, n$ with $m<n$ 
\item not differentiable for each pair $m, n$ with $m>n$ 
\item not differentiable for each pair $m, n$ with $m<n$

\end{enumerate}
\hfill [JAM 2015]




\item  

Let $G$ and $H$ be nonempty subsets of $\mathbb{R}$, where $G$ is connected and $G \cup H$ is not connected. Which one of the following statements is true for all such $G$ and $H$ ? 
\begin{enumerate}
    \item If $G \cap H=\emptyset$, then $H$ is connected 
\item If $G \cap H=\emptyset$, then $H$ is not connected 
\item If $G \cap H \neq \emptyset$, then $H$ is connected 
\item If $G \cap H \neq \emptyset$, then $H$ is not connected
\end{enumerate}
\hfill [JAM 2015]





\item  

Let $f:\left\{(x, y) \in \mathbb{R}^{2}: x>0, y>0\right\} \rightarrow \mathbb{R}$ be given by
$$
f(x, y)=x^{\frac{1}{3}} y^{\frac{-4}{3}} \tan ^{-1}\left(\frac{y}{x}\right)+\frac{1}{\sqrt{x^{2}+y^{2}}}
$$
Then the value of

$$
g(x, y)=\frac{x f_{x}(x, y)+y f_{y}(x, y)}{f(x, y)}
$$

\begin{enumerate}
    \item changes with $x$ but not with $y$ 
\item changes with $y$ but not with $x$ 
\item changes with $x$ and also with $y$ 
\item neither changes with $x$ nor with $y$
\end{enumerate}
\hfill [JAM 2015]







Q.1

Let $f: \mathbb{R} \rightarrow \mathbb{R}$ be a function defined by $f(x)=\int_{-5}^{x}(t-1)^{3} d t$.

In which of the following interval(s), $f$ takes the value 1 ?\\

\begin{enumerate}
    
\item $[-6,0]$\\
\item $[-2,4]$\\
\item $[2,8]$\\
\item $[6,12]$
\end{enumerate}
\hfill [JAM 2015]




\item  Which of the following statements is (are) true?\\

\begin{enumerate}
    \item $\mathrm{z}_{2} \oplus \mathrm{z}_{3}$ is isomorphic to $\mathrm{Z}_{6}$\\
\item $\mathrm{Z}_{3} \oplus \mathrm{Z}_{3}$ is isomorphic to $\mathrm{Z}_{9}$\\
\item $\mathrm{Z}_{4} \oplus \mathrm{Z}_{6}$ is isomorphic to $\mathrm{Z}_{24}$\\
\item $\mathrm{z}_{2} \oplus \mathrm{z}_{3} \oplus \mathrm{z}_{5}$ is isomorphic to $\mathrm{z}_{30}$
\end{enumerate}
\hfill [JAM 2015]



 Which of the following conditions implies (imply) the convergence of a sequence $\left\{x_{n}\right\}$ of real numbers?
\begin{enumerate}
    \item Given $\varepsilon>0$ there exists an $n_{0} \in \mathbb{N}$ such that for all $n \geq n_{0},\left|x_{n+1}-x_{n}\right|<\varepsilon$
\item Given $\varepsilon>0$ there exists an $n_{0} \in \mathbb{N}$ such that for all $n \geq n_{0}, \frac{1}{(n+1)^{2}}\left|x_{n+1}-x_{n}\right|<\varepsilon$
\item Given $\varepsilon>0$ there exists an $n_{0} \in \mathbb{N}$ such that for all $n \geq n_{0},(n+1)^{2}\left|x_{n+1}-x_{n}\right|<\varepsilon$
\item Given $\varepsilon>0$ there exists an $n_{0} \in \mathbb{N}$ such that for all $m, n$ with $m>n \geq n_{0},\left|x_{m}-x_{n}\right|<\varepsilon$
\end{enumerate}
\hfill [JAM 2015]





\item  

Let $\vec{F}$ be a vector field given by $\vec{F}(x, y, z)=-y \hat{\imath}+2 x y \hat{\jmath}+z^{3} \hat{k}$, for $(x, y, z) \in \mathbb{R}^{3}$. If $C$ is the curve of intersection of the surfaces $x^{2}+y^{2}=1$ and $y+z=2$, then which of the following is (are) equal to $\left|\int_{C} \vec{F} \cdot d \vec{r}\right| ?$

\begin{enumerate}
    
\item $\int_{0}^{2 \pi} \int_{0}^{1}(1+2 r \sin \theta) r d r d \theta$\\
\item $\int_{0}^{2 \pi}\left(\frac{1}{2}+\frac{2}{3} \sin \theta\right) d \theta$\\
\item $\int_{0}^{2 \pi} \int_{0}^{1}(1+2 r \sin \theta) d r d \theta$\\
\item $\int_{0}^{2 \pi}(1+\sin \theta) d \theta$
\end{enumerate}
\hfill [JAM 2015]




\item  

Let $V$ be the set of $2 \times 2$ matrices $\left[\begin{array}{ll}a_{11} & a_{12} \\ a_{21} & a_{22}\end{array}\right]$ with complex entries such that $a_{11}+a_{22}=0$. Let $W$ be the set of matrices in $V$ with $a_{12}+\overline{a_{21}}=0$. Then, under usual matrix addition and scalar multiplication, which of the following is (are) true?

\begin{enumerate}
    \item $V$ is a vector space over $\mathbb{C}$
\item $W$ is a vector space over $\mathbb{C}$
\item $V$ is a vector space over $\mathbb{R}$
\item $W$ is a vector space over $\mathbb{R}$

\end{enumerate}
\hfill [JAM 2015]




\item  

The initial value problem
$$
y^{\prime}=\sqrt{y}, y(0)=\alpha, \quad \alpha \geq 0
$$
has

\begin{enumerate}
    \item at least two solutions if $\alpha=0$
\item no solution if $\alpha>0$
\item at least one solution if $\alpha>0$
\item a unique solution if $\alpha=0$
\end{enumerate}
\hfill [JAM 2015]





\item  Which of the following statements is (are) true on the interval $\left(0, \frac{\pi}{2}\right)$ ?

\begin{enumerate}
    \item $\cos x<\cos (\sin x)$
\item $\tan x<x$
\item $\sqrt{1+x}<1+\frac{x}{2}-\frac{x^{2}}{8}$
\item $\frac{1-x^{2}}{2}<\ln (2+x)$
\end{enumerate}
\hfill [JAM 2015]





\item  Let $f: \mathbb{R}^{2} \rightarrow \mathbb{R}$ be defined by
$$
f(x, y)=\left\{\begin{array}{cc}
x y \frac{x^{2}-y^{2}}{x^{2}+y^{2}}, & (x, y) \neq(0,0) \\
0, & (x, y)=(0,0)
\end{array}\right.
$$
At $(0,0)$,\\

\begin{enumerate}
    \item $f$ is not continuous
\item $f$ is continuous, and both $f_{x}$ and $f_{y}$ exist
\item $f$ is differentiable
\item $f_{x}$ and $f_{y}$ exist but $f$ is not differentiable
\end{enumerate}
\hfill [JAM 2015]



 \item  Let $f, g:[0,1] \rightarrow[0,1]$ be functions. Let $R(f)$ and $R(g)$ be the ranges of $f$ and $g$, respectively. Which of the following statements is (are) true?
 
\begin{enumerate}
    

\item If $f(x) \leq g(x)$ for all $x \in[0,1]$, then $\sup R(f) \leq \inf R(g)$
\item If $f(x) \leq g(x)$ for some $x \in[0,1]$, then inf $R(f) \leq \sup R(g)$
\item If $f(x) \leq g(y)$ for some $x, y \in[0,1]$, then inf $R(f) \leq \sup R(g)$
\item If $f(x) \leq g(y)$ for all $x, y \in[0,1]$, then sup $R(f) \leq \inf R(g)$

\end{enumerate}
\hfill [JAM 2015]


\item  Let $f:(-1,1) \rightarrow \mathbb{R}$ be the function defined by
$$
f(x)=x^{2} e^{1 /\left(1-x^{2}\right)}
$$
Then
\begin{enumerate}
    
\item $f$ is decreasing in $(-1,0)$
\item $f$ is increasing in $(0,1)$
\item $f(x)=1$ has two solutions in $(-1,1)$
\item $f(x)=1$ has no solutions in $(-1,1)$
\end{enumerate}
\hfill [JAM 2015]




\item  Let $C$ be the straight line segment from $P(0, \pi)$ to $Q\left(4, \frac{\pi}{2}\right)$, in the $x y$-plane. Then the value of $\int_{C} e^{x}(\cos y d x-\sin y d y)$ is \rule{3cm}{1pt}.
\hfill [JAM 2015]




\item  Let $S$ be the portion of the surface $z=\sqrt{16-x^{2}}$ bounded by the planes $x=0, x=2, y=0$, and $y=3$. The surface area of $S$, correct up to three decimal places, is \rule{3cm}{1pt}.
\hfill [JAM 2015]




\item  The number of distinct normal subgroups of $S_{3}$ is \rule{3cm}{1pt}.
\hfill [JAM 2015]




\item  Let $f: \mathbb{R}^{2} \rightarrow \mathbb{R}$ be defined by
$$
f(x, y)=\left\{\begin{aligned}
\left(1+\frac{x}{y}\right), & y \neq 0 \\
y &=0
\end{aligned}\right.
$$
If the directional derivative of $f$ at $(0,0)$ exists along the direction $\cos \alpha \hat{\imath}+\sin \alpha \hat{j}$, where $\sin \alpha \neq 0$, then the value of $\cot \alpha$ is \rule{3cm}{1pt}.
\hfill [JAM 2015]



\item  Let $f: \mathbb{R}^{3} \rightarrow \mathbb{R}$ be defined by
$$
f(x, y, z)=\sin x+2 e^{\frac{y}{2}}+z^{2}
$$
The maximum rate of change of $f$ at $\left(\frac{\pi}{4}, 0,1\right)$, correct up to three decimal places, is \rule{3cm}{1pt}.
\hfill [JAM 2015]



\item  

If the power series
$$
\sum_{n=0}^{\infty} \frac{n !}{n^{n}} x^{2 n}
$$
converges for $|x|<c$ and diverges for $|x|>c$, then the value of $c$, correct up to three decimal places, is \rule{3cm}{1pt}.
\hfill [JAM 2015]




\item  If $5^{2015}=n(\bmod 11)$ and $n \in\{0,1,2,3,4,5,6,7,8,9,10\}$, then $n$ is equal to \rule{3cm}{1pt}.
\hfill [JAM 2015]




\item  If the set $\left\{\left[\begin{array}{cc}x & -x \\ -1 & 0\end{array}\right],\left[\begin{array}{cc}0 & -1 \\ x & x\end{array}\right],\left[\begin{array}{cc}1 & -1 \\ 1 & 0\end{array}\right]\right\}$ is linearly dependent in the vector space of all $2 \times 2$ matrices with real entries, then $x$ is equal to Let $f: \mathbb{R} \rightarrow \mathbb{R}$ be defined by
$$
f(x)= \begin{cases}x^{6}-1, & x \in \mathbb{Q} \\ 1-x^{6}, & x \notin \mathbb{Q}\end{cases}
$$
The number of points at which $f$ is continuous, is \rule{3cm}{1pt}.
\hfill [JAM 2015]



\item  

Let $f:(0,1) \rightarrow \mathbb{R}$ be a continuously differentiable function such that $f^{\prime}$ has finitely many zeros in $(0,1)$ and $f^{\prime}$ changes sign at exactly two of these points. Then for any $y \in \mathbb{R}$, the maximum number of solutions to $f(x)=y$ in $(0,1)$ is \rule{3cm}{1pt}.
\hfill [JAM 2015]



Q. $11-$ Q. 20 carry two marks each.

\item  Let $R$ be the planar region bounded by the lines $x=0, y=0$ and the curve $x^{2}+y^{2}=4$, in the first quadrant. Let $C$ be the boundary of $R$, oriented counter-clockwise. Then the value of
$$
\oint_{C} x(1-y) d x+\left(x^{2}-y^{2}\right) d y
$$
is \rule{3cm}{1pt}.
\hfill [JAM 2015]




\item  Suppose $G$ is a cyclic group and $\sigma, \tau \in G$ are such that $\operatorname{order}(\sigma)=12$ and $\operatorname{order}(\tau)=21$. Then the order of the smallest group containing $\sigma$ and $\tau$ is \rule{3cm}{1pt}.
\hfill [JAM 2015]




\item  The limit

$$
\lim _{n \rightarrow \infty} \sum_{k=2}^{n} \frac{1}{k^{3}-k}
$$
is equal to \rule{3cm}{1pt}.
\hfill [JAM 2015]





\item  Let $M_{2}(\mathbb{R})$ be the vector space of $2 \times 2$ real matrices. Let $V$ be a subspace of $M_{2}(\mathbb{R})$ defined by
$$
V=\left\{A \in M_{2}(\mathbb{R}): A\left[\begin{array}{ll}
0 & 2 \\
3 & 1
\end{array}\right]=\left[\begin{array}{ll}
0 & 2 \\
3 & 1
\end{array}\right] A\right\}
$$
Then the dimension of $V$ is \rule{3cm}{1pt}.
\hfill [JAM 2015]




\item  

Let $f: \mathbb{R}^{2} \rightarrow \mathbb{R}$ be defined by

Then the integral
$$
f(x, y)=\left\{\begin{array}{cc}
\frac{y}{\sin y}, & y \neq 0 \\
1, & y=0
\end{array}\right.
$$
$$
\frac{1}{\pi^{2}} \int_{x=0}^{1} \int_{y=\sin ^{-1} x}^{\frac{\pi}{2}} f(x, y) d y d x
$$
correct up to three decimal places \rule{3cm}{1pt}.
\hfill [JAM 2015]




\item  

The coefficient of $\left(x-\frac{\pi}{4}\right)^{3}$ in the Taylor series expansion of the function
$$
f(x)=3 \sin x \cos \left(x+\frac{\pi}{4}\right), \quad x \in \mathbb{R}
$$
about the point $\frac{\pi}{4}$, correct up to three decimal places, is \rule{3cm}{1pt}.
\hfill [JAM 2015]



\item  

If $\int_{0}^{x}\left(e^{-t^{2}}+\cos t\right) d t$ has the power series expansion $\sum_{n=1}^{\infty} a_{n} x^{n}$, then $a_{5}$, correct up to three decimal places, is equal to \rule{3cm}{1pt}.
\hfill [JAM 2015]



\item  

Let $\ell$ be the length of the portion of the curve $x=x(y)$ between the lines $y=1$ and $y=3$, where $x(y)$ satisfies
$$
\frac{d x}{d y}=\frac{\sqrt{1+y^{2}+y^{4}}}{y}, \quad x(1)=0
$$
The value of $\ell$, correct upto three decimal places, is
\hfill [JAM 2015]

\item 
The limit
$$
\lim _{x \rightarrow 0+} \frac{9}{x}\left(\frac{1}{\tan ^{-1} x}-\frac{1}{x}\right)
$$
is equal to \rule{3cm}{1pt}.
\hfill [JAM 2015]




\item  Let $P$ and $Q$ be two real matrices of size $4 \times 6$ and $5 \times 4$, respectively. If rank $(Q)=4$ and $\operatorname{rank}(Q P)=2$, then $\operatorname{rank}(P)$ is equal to \rule{3cm}{1pt}.
\hfill [JAM 2015]



   
\end{enumerate}
\newpage
